\documentclass[11pt, a4paper, german]{article}

\usepackage[utf8]{inputenc}
\usepackage[ngerman]{babel}
\usepackage{amsmath,amsthm,amssymb}
\usepackage[table]{xcolor} % Für farbige Tabellenzellen
\usepackage{hhline} % mehr Kontrolle über Tabellenlinien
\usepackage{tikz}
\usepackage{enumitem} % bessere Aufzählungen
\usepackage{mathtools} % \coloneqq
\usepackage{stmaryrd} % \mapsfrom
\usetikzlibrary{cd}
\usepackage{geometry}
\geometry{margin=2.5cm}
\usepackage{xfrac}
%\usepackage{color} % \red
%\usepackage{wrapfig} % Bild neben Text
\usepackage{cutwin} % Bild neben Text
\usepackage{mathabx} % \divides
\usepackage{array} % für \setlength\extrarowheight{2pt}
\usepackage{lscape} % Für Querformat

\usepackage[style=alphabetic-verb, backend=bibtex]{biblatex} % Literaturverzeichnis
\bibliography{literatur.bib}
%\usepackage[natbib=true,style=numeric]{biblatex}
%\bibliography{bibliography}

\theoremstyle{definition}
%\newtheorem*{nota}{Notation}
\newtheorem{bsp}{Beispiel}
\newtheorem{satz}{Satz}
\newtheorem{lem}{Lemma}
\newtheorem{defn}{Definition}
\newtheorem{beob}{Beobachtung}
\newtheorem{bspe}{Beispiele}
\newtheorem{kor}{Korollar}

\theoremstyle{remark}
\newtheorem*{bem}{Bemerkung}

\newcommand{\TODO}[1]{\textcolor{red}{TODO: #1}} %\red{#1}} % TODO-Markierungen

% Zahlbereiche
\newcommand{\R}{\mathbb{R}} % Reelle Zahlen
\newcommand{\N}{\mathbb{N}} % Natürliche Zahlen
\newcommand{\Z}{\mathbb{Z}} % Ganze Zahlen
%\newcommand{\C}{\mathbb{C}} % Komplexe Zahlen
\newcommand{\Q}{\mathbb{Q}} % Rationale Zahlen
\newcommand{\ZZ}[2]{\mathbb{Z}_{\text{\scriptsize \,$#1$}}^{\text{\scriptsize \,$#2$}}} % Ganze Zahlen mod n hoch m

% Schöne Fürall- und Existenzquantoren
\newcommand{\fa}[1]{\forall \, {#1} \,:\,}
%\newcommand{\ex}[1]{\exists \, {#1} \,:\,}

\DeclareMathOperator{\Hom}{Hom} % Homomorphisms
\DeclareMathOperator{\id}{id} % Identität
\DeclareMathOperator{\spann}{spann} % Spann
\DeclareMathOperator{\im}{im} % Image (Bild)
\DeclareMathOperator{\coker}{coker} % Kokern
\newcommand{\pt}{\mathrm{pt}}
\newcommand{\blank}{\text{--}} % Platzhalter
\newcommand{\ntimes}{\!\times\!} % schmaleres (narrower) \times
\newcommand{\angles}[1]{{\langle #1 \rangle}}
\newcommand{\const}[1]{\text{konst } #1} % konstante Funktion mit Wert #1
\newcommand{\LH}{\mathcal{H}} % Local homology
\DeclareMathOperator{\colim}{colim} % Kolimes
\newcommand{\nspace}[1]{\foreach \i in {1,...,#1}{ \! }} % Negativer Abstand
\DeclareMathOperator{\Tor}{Tor} % Tor-Funktor
\newcommand{\SC}{\mathcal{C}} % Serre-Klasse
\newcommand{\FG}{\mathcal{FG}} % Serre-Klasse der endlich erzeugten Gruppen
\newcommand{\F}{\mathcal{F}} % Serre-Klasse der endlichen Gruppen

% Abkürzungen
\newcommand{\ES}{Es sei} % Prof. Hanke schreibt das immer, mag das anscheinend lieber als ein bloßes "Sei"
\newcommand{\ESn}{Es seien} % Prof. Hanke schreibt das immer, mag das anscheinend lieber als ein bloßes "Sei"
\newcommand{\zB}{z.\,B.}
\renewcommand{\dh}{d.\,h.} % das heißt

% Intervalle
\newcommand{\cinterval}[2]{\left[ #1, #2 \right]} % closed interval
\newcommand{\I}{I} % Kompaktes Einheitsinterval [0,1]
%\newcommand{\I}{\cinterval{0}{1}} % Kompaktes Einheitsinterval [0,1] (alternativ)

\newcommand{\fall}[1]{\emph{#1}:\enspace}

% Schöne Mengen { #1 | #2 } (benötigt mathtools)
% siehe http://tex.stackexchange.com/questions/13634/define-pretty-sets-in-latex-esp-how-to-do-the-condition-separator
\DeclarePairedDelimiterX\Set[2]{\lbrace}{\rbrace}%
 { #1 \,\delimsize|\, #2 }

% http://tex.stackexchange.com/questions/117732/tikz-and-babel-error
% Es ist schierer Wahnsinn, welche Hacks LaTeX benötigt!
\tikzset{
  every picture/.prefix style={
    execute at begin picture=\shorthandoff{"}
  }
}

% Zentrierte kommutative Diagramme
\newenvironment{centertikzcd}
  {\begin{center}\begin{tikzcd}}
  {\end{tikzcd}\end{center}}

% TikZ-Makros
\newcommand{\zeroDot}[2]{\node[draw,circle,inner sep=0.4pt,fill] at (#1,#2) {};} % Punkt, der in Spektralsequenz für eine Null steht

\begin{document}


\begin{titlepage}
  \begin{center}
    \ \vspace{0.1cm}\\
    \Huge \textbf{\textsc{Der Endlichkeitssatz von Serre über die Homotopiegruppen der Sphären}} \vspace{0.2cm}\\
    \vspace{2cm}
    \LARGE \textbf{\textsc{Bachelorarbeit}}\\
    von\\
    \textbf{Tim Baumann}\\
    \vspace{1.1cm}
    \begin{figure}[htb]
      \centering
      %\includegraphics [scale=0.19] {D:/Studium_Mathe_Augsburg_6/Clifford-Bundel/SignUA.pdf}
    \end{figure}
    \vspace{1.1cm}
    \large Eingereicht am \\
    22. 12. 2015 \\[1cm]
    \includegraphics[scale=0.4]{uni_aux_logo.jpg} \\[1cm]
    %\large Institut für Mathematik\\ Universität Augsburg \\
    \vspace{1.2cm}
    \large \begin{tabular}{ll}
      Erstgutachter: & Prof. Dr. Bernhard Hanke \\
      Zweitgutachter: & \TODO{???}
    \end{tabular}
  \end{center}
\end{titlepage}


\tableofcontents
\newpage

\section{Einleitung}

% Hatcher, 1.21
\begin{satz}
  Die Homotopiegruppen $\pi_i(S^n, *)$,~$i > n$, sind endlich bis auf die Gruppen $\pi_{4k-1}(S^{2k})$, $k \geq 1$, welche jeweils isomorph zu einer direkten Summe von $\Z$ und einer endlichen Gruppe sind.
\end{satz}

\section{Homotopiegruppen der Sphären}

\begin{defn}
  Es sei $(X, x_0)$ ein punktierter topologischer Raum.
  Für $n \geq 0$ sei
  \[ \pi_n(X, x_0) \coloneqq [(S^n, *), (X, x_0)] \]
  die Menge der basispunkterhaltenden Abbildungen $S^n \to X$ modulo basispunkterhaltender Homotopie.
  Für $n \geq 1$ heißt $\pi_n(X)$ die \emph{$n$-te Homotopiegruppe} von $X$ (mit Basispunkt $x_0$).
\end{defn}

Die Gruppenstruktur auf $\pi_n(X)$ wird induziert durch die Kogruppenstruktur auf $S^n$, welche durch die Abbildung $S^n \to S^n \vee S^n$ gegeben ist, die die Südkappe $D^n_+$ auf die erste $S^n$, die Nordkappe $D^n_{-}$ auf die zweite $S^n$ und den Äquator $D^n_{+} \cap D^n_{-}$ auf den Anklebepunkt von $S^n \vee S^n$ abbildet.
Das neutrale Element der Homotopiegruppen wird repräsentiert durch die konstante Abbildung auf den Basispunkt $x_0$.
Man kann zeigen, dass $\pi_n(X, x_0)$ für $n \geq 2$ abelsch ist (vgl. \cite{hatcher:at}).

Eine basispunkterhaltende Abbildung $g : (X, x_0) \to (Y, y_0)$ induziert Abbildungen $g_* = \pi_i(g) : \pi_i(X, x_0) \to \pi_i(Y, y_0)$ durch Nachkomponieren mit $g$.
Auf diese Weise wird $\pi_i$ zu einem Funktor von der Kategorie der punktierten topologischen Räume in die Kategorie der (abelschen) Gruppen.

Die Menge $\pi_0(X, x_0)$ ist die Menge der Wegzusammenhangskomponenten von $X$.
Wir schreiben $\pi_0(X, x_0) = 0$, falls die einzige Wegzusammenhangskomponente von $X$ die von $x_0$ ist.

Falls ein Weg von $x_0$ nach $x'_0$ in $X$ existiert, so sind $\pi_n(X, x_0)$ und $\pi_n(X, x'_0)$ isomorph.
Für nicht leere, wegzusammenhängende Räume $X$ kann man daher den Basispunkt in der Notation weglassen von der Homotopiegruppe $\pi_n(X)$ von $X$ sprechen.

% TODO: Motivation für Studium der Homotopiegruppen der Sphären

Für $i < n$ gilt $\pi_i(S^n) = 0$:
Wir können die Sphären $S^i$ und $S^n$ als CW-Komplexe mit jeweils einer Nullzelle und einer $i$-Zelle bzw. $n$-Zelle realisieren.
Das zelluläre Approximationstheorem besagt, dass jede Abbildung $f : S^i \to S^n$ homotop relativ Basispunkt zu einer zellulären Abbildung ist, \dh{} einer Abbildung $\tilde{f} : S^i \to S^n$, die das $i$-Skelett von $S^i$ (das ist ganz $S^i$) auf das $i$-Skelett von $S^n$ (das ist $\{ * \}$) abbildet.
In anderen Worten ist jede Abbildung $S^i \to S^n$ homotop zur konstanten Abbildung und repräsentiert daher das neutrale Element in $\pi_i(S^n)$.
% TODO: Hatcher zitieren

Mit der universellen Überlagerung $p : \R \to S^1, \enspace t \mapsto e^{it}$ kann man die Homotopiegruppen von $S^1$ bestimmen:
Die Fundamentalgruppe von $S^1$ ist isomorph zur Decktransformationsgruppe dieser Überlagerung, von der man leicht zeigen kann, dass diese isomorph zu $\Z$ ist.
Es sei nun $i > 1$ und $f : S^i \to S^1$ stetig.
Dann existiert eine Hochhebung $\tilde{f} : S^i \to \R$, da $f_* (\pi_1(S^n)) = f_*(0) = 0 = p_*(\pi_0(\R))$.
Da $\R$ zusammenziehbar ist, finden wir eine Homotopie $H$ von $\tilde{f}$ zur konstanten Abbildung $x \mapsto 0$.
Dann ist $p \circ H$ eine Homotopie zwischen $f$ und der konstanten Abbildung.
Somit ist $\pi_i(S^1) = 0$ für $i > 1$.

\begin{defn}
  Die \emph{(reduzierte) Einhängung} $\Sigma X$ eines punktierten Raumes $(X, x_0)$ ist
  \[ \Sigma X \coloneqq (X \times I) / (X \times \{ 0, 1 \} \cup \{ x_0 \} \times I). \]
  Der Basispunkt $*$ von $\Sigma X$ ist der auf einen Punkt zusammengezogene Teilraum.
\end{defn}

Einhängung ist ein Endofunktor der Kategorie der punktierten topologischen Räume: \\
Für $f : (X, x_0) \to (Y, y_0)$ ist
\[
  \Sigma f : (\Sigma X, *) \to (\Sigma Y, *), \quad
  [(x, t)] \mapsto [(f(x), t)].
\]

Die Einhängung eines CW-Komplexes ist wieder ein CW-Komplex.
Man sieht leicht, dass $\Sigma S^i \approx S^{i+1}$.
Somit induziert Einhängung eine Abbildung $E : \pi_i(X, x_0) \to \pi_i(\Sigma X, *)$.
Da Einhängung mit den Kogruppenstruktur von $S^i$ und $S^{i+1}$ verträglich ist, ist $E$ sogar ein Gruppenhomomorphismus.

\begin{defn}
  Ein nicht leerer topologischer Raum $X$ heißt \emph{$n$-zusammenhängend}, falls $\pi_i(X) = 0$ für $0 \leq i \leq n$.
\end{defn}

% TODO: Hatcher zitieren, p. 360
\begin{satz}[Freudenthal'scher Einhängungssatz]
  Es sei $n \geq 0$ und $X$ ein $n$-zusammenhängender CW-Komplex.
  Dann ist $E : \pi_i(X) \to \pi_{i+1}(\Sigma X)$ bijektiv für $0 \leq i \leq 2n$ und surjektiv für $i = 2n + 1$.
\end{satz}

Insbesondere ist die $i$-fach iterierte Einhängung $\Sigma^i X$ eines $n$-zusammenhängenden CW-Komplexes $X$ ein $(n{+}i)$-zusammenhängender CW-Komplex.

Für einen beliebigen CW-Komplex $X$ gilt daher $\pi_{n+j}(\Sigma^j X) \cong \pi_{n+j+1}(\Sigma^{j+1} X)$ für $n + j \leq 2(j-1) \Leftrightarrow n+2 \leq j$.
Somit ist
\[ \pi_{2n+2}(\Sigma^{n} X) \cong \pi_{2n+3}(\Sigma^{n+1} X) \cong \pi_{2n+4}(\Sigma^{n+2} X) \cong \ldots \cong \colim \pi_{n+j}(\Sigma^j X). \]
Diese Gruppe heißt \emph{$n$-te stabile Homotopiegruppe} von $X$.
Im Fall der Sphären haben wir mit $X = S^0$:
\[ \pi_{2n+2}(S^{n+2}) \cong \pi_{2n+3}(S^{n+3}) \cong \pi_{2n+4}(S^{n+4}) \cong \ldots \cong \colim_j \pi_{n+j}(S^j) =: \pi^s_n. \]
Wir behaupten, dass $E : \pi_1(S^1) \to \pi_2(S^2)$ ein Isomorphismus ist.
Aus dem Einhängungssatz folgt, dass diese Abbildung surjektiv ist.
Somit ist $\pi_2(S^2)$ ein zyklische Gruppe mit Erzeuger $\id_{S^2}$.
Es sei nun $k \in \Z$.
Dann ist $H_2(k \cdot \id_{S^2}) : H_2(S^2) \to H^2(S^2)$ gegeben durch Multiplikation mit $k$.
Angenommen, $k \cdot \id_{S^2}$ ist nullhomomotop.
Dann ist auch $H_2(k \cdot \id_{S^2}) = 0$ und somit $k = 0$.
Somit ist $E$ injektiv.
Es folgt
\[ \Z \cong \pi_1(S^1) \cong \pi_2(S^2) \cong \ldots \cong \pi_i(S^i)  \cong \ldots \cong \colim_j \pi_j(S^j) = \pi^s_0. \]
Dabei ist $\id_{S^i}$ ein Erzeuger von $\pi_i(S^i)$.

Zusammengefasst wissen wir also
\[
  \pi_i(S^n, *) \cong \begin{cases}
    0, & \text{für $i < n$ und $i > n = 1$} \\
    \Z, & \text{für $i = n$,} \\
    \Z \oplus \text{endliche Gruppe}, & \text{für $i=2n-1$ und $n$ gerade,} \\
    \text{endliche Gruppe}, & \text{sonst.}
  \end{cases}
\]

\section{Faserungen}

\begin{defn}
  Eine \emph{Serre-Faserung} ist eine stetige Abbildung $p : E \to B$, welche die \emph{Homotopieliftungseigenschaft} (HLE) für die Scheiben $D^n$ besitzt, \dh{}
  für alle $n \geq 0$ und für alle stetigen Abbildungen $H$, $H_0$ wie unten, sodass das äußere Quadrat kommutiert, gibt es eine stetige Abbildung $\tilde{H}$, sodass die beiden Dreiecke kommutieren:
  \begin{centertikzcd}[row sep=1.2cm, column sep=1.4cm]
    D^n \arrow[r, "H_0"] \arrow[d, hook, "i_0", swap] &
    E \arrow[d, "p"] \\
    D^n \times \I \arrow[r, "H"] \arrow[ur, "\exists\, \tilde{H}", dashed] &
    B
  \end{centertikzcd}
  Dabei ist $i_0$ die Inklusion von $D^n$ in $D^n \times \I$ als $D^n \times \{ 0 \}$. \\
  Eindeutigkeit von $\tilde{H}$ wird nicht gefordert.
\end{defn}

\begin{lem}
  \ES{} $p : E \to B$ eine stetige Abbildung. Dann sind äquivalent:
  \begin{enumerate}[label=\alph*)]
    \item $p$ ist eine Serre-Faserung
    \item $p$ besitzt die \emph{relative Homotopieliftungseigenschaft} für CW-Paare, \dh{} für alle CW-Paare $(X, A)$ und für alle $H_0$ und $H$ wie unten, sodass das äußere Quadrat kommutiert, gibt eine stetige Abbildung $\tilde{H}$, sodass die beiden Dreiecke kommutieren:
    \begin{centertikzcd}[row sep=1.2cm, column sep=1.4cm]
      X \!\times\! \{ 0 \} \cup A \!\times\! \I \arrow[r, "H_0"] \arrow[d, hook, swap] &
      E \arrow[d, "p"] \\
      X \times \I \arrow[r, "H"] \arrow[ur, "\exists\, \tilde{H}", dashed] &
      B
    \end{centertikzcd}
  \end{enumerate}
\end{lem}

\begin{bem}
  Eine \emph{Hurewicz-Faserung} ist eine Serre-Faserung, welche die Homotopieliftungseigenschaft sogar für alle topologischen Räume besitzt.
\end{bem}

% vgl. Diskussion im Hatcher auf Seite 376
\begin{proof}
  "`b) $\implies$ a)"' \enspace Folgt sofort mit $(X, A) \coloneqq (D^n, \emptyset)$. \\[2pt]
  "`a) $\implies$ b)"' \enspace Wir behandeln zunächst den Fall $(X, A) = (D^n, S^{n-1})$, $n \in \N$. Dann ist $(D^n \times \I, D^n \times \{ 0 \} \cup S^{n-1} \cup \I) \approx (D^n)$ homöomorph als Raumpaar.
  Somit ist die relative Homotopieliftungseigenschaft in diesem Fall gleichbedeutend zur Homotopieliftungseigenschaft für die Scheibe $D^n$. \\
  \ES{} nun $(X, A)$ ein beliebiges Raumpaar.
  % TODO: Ausführlicher?
  Dann kann man induktiv die Homotopie $H$ auf die $i$-Zellen $e^i_\alpha$ von $X \setminus A$ fortsetzen.
  Dabei ist die Homotopie auf $S^{n-1} = \partial D^n$ durch die Komposition der bisher konstruierten Homotopie mit der anheftenden Abbildung $\phi_\alpha : S^{n-1} \to X^{n-1}$ vorgegeben.
  Man erhält die Fortsetzung durch Anwenden des zuerst bewiesenen Falls.
\end{proof}

% Lemma 4.5 in http://www.math.washington.edu/~mitchell/Notes/serre.pdf
\begin{lem}
  \ESn{} $p : E \to B$ eine Serre-Faserung, $b_0 \in B$, $F \coloneqq p^{-1}(b_0)$ die Faser über $b_0$ und $f_0 \in F$.
  Dann gibt es eine lange exakte Sequenz
  \[ \ldots \to \pi_n(F, f_0) \xrightarrow{i_*} \pi_n(E, f_0) \xrightarrow{p_*} \pi_n(B, b_0) \xrightarrow{\partial} \pi_{n-1}(F, f_0) \to \ldots \to \pi_1(B, b_0) \]
  von Homotopiegruppen.
  Dabei ist $i : F \hookrightarrow E$ die Inklusion.
\end{lem}

\begin{proof}
  Die gesuchte exakte Sequenz ist die lange exakte Homotopiesequenz
  \[ \ldots \to \pi_n(F, f_0) \xrightarrow{i_*} \pi_n(E, f_0) \to \pi_n(E, F, f_0) \xrightarrow{\partial} \pi_{n-1}(F, f_0) \to \ldots \to \pi_1(E, F, f_0) \]
  des Raumpaares $(E, F)$.
  Es bleibt zu zeigen: $\pi_n(E, F, f_0) \cong \pi_n(B, b_0)$ als Gruppe für $n > 1$ und als punktierte Menge für $n = 1$.
  Der Isomorphismus muss außerdem so gewählt werden, dass
  \[ p_* = \left( \pi_n(E, f_0) \to \pi_n(E, F, f_0) \xrightarrow{\cong} \pi_n(B, b_0) \right). \]
  Wir zeigen: $p_* : \pi_n(E, F, f_0) \to \pi_n(B, b_0)$ ist der gesuchte Isomorphismus (damit ist obige Gleichung erfüllt). \\[2pt]
  \emph{Surjektivität}: Sei $[g : (I^{n+1}, \partial I^{n+1}, b_0) \to (B, \{ b_0 \}, b_0)] \in \pi_{n+1}(B, b_0)$, $n \geq 0$.
  Sei $\tilde{g}$ der Lift im folgenden relativen HLE-Diagramm:
  \begin{centertikzcd}[row sep=1.2cm, column sep=1.4cm]
    U \arrow[r, "\const{f_0}"] \arrow[d, hook, swap] &
    E \arrow[d, "p"] \\
    I^n \times I \arrow[r, "g"] \arrow[ur, "\exists\, \tilde{g}", dashed] \arrow[r, "g"] &
    B
  \end{centertikzcd}
  wobei $U \coloneqq I^n \!\times\! \{ 0 \} \cup (\partial I^n) \!\times\! I \subset I^{n+1}$.
  Dann kann man $\tilde{g}$ als eine Abbildung $(I^{n+1}, \partial I^{n+1}, U) \to (E, F, \{ f_0 \})$ von Raumtripeln auffassen, welche ein Element von $\pi_{n+1}(E, F, f_0)$ repräsentiert.
  Es gilt $p_*[\tilde{g}] = [p \circ \tilde{g}] = [g]$. \\[2pt]
  \emph{Injektivität}: Seien $[h_0], [h_1] \in \pi_{n+1}(E, F, f_0)$ mit $p_*[h_0] = p_*[h_1]$.
  Sei
  \[ H : I \times I^{n+1}, \quad (t, x) \mapsto H_t(x) \]
  eine Homotopie mit $H_0 = p \circ h_0$, $H_1 = p \circ h_1$, welche zu jedem Zeitpunkt $t \in \I$ eine Abbildung $H_t : (I^{n+1}, \partial I^{n+1}) \to (B, \{ b_0 \})$ von Raumpaaren ist.
  Betrachte folgendes HLE-Diagramm:
  \begin{centertikzcd}[row sep=1.2cm, column sep=1.8cm]
    V \arrow[r, "h"] \arrow[d, hook, swap] &
    E \arrow[d, "p"] \\
    I^{n+1} \times I \arrow[r, "H"] \arrow[ur, "\exists\, \tilde{H}", dashed] &
    B
  \end{centertikzcd}
  mit $V \coloneqq I^{n+1} \!\times\! \{ 0 \} \cup (\partial I^{n+1}) \!\times\! I \subset I^{n+2}$ und
  \[
    h|_{\{0\} \times I^{n+1}} \coloneqq h_0, \quad
    h|_{\{1\} \times I^{n+1}} \coloneqq h_1, \quad
    h|_{I \times U} \coloneqq \const{f_0}.
  \]
  Nun ist $\tilde{H}$ eine Homotopie von $h_0$ nach $h_1$, welche zu jedem Zeitpunkt $t$ eine Abbildung $\tilde{H}_t : (I^{n+1}, \partial I^{n+1}, U) \to (E, F, \{ b_0 \})$ von Raumtripeln ist.
\end{proof}

\begin{defn}
  \ESn{} $p : E \to B$ und $g : X \to B$ stetig.
  Der \emph{Pullback} von $p$ entlang $g$ ist die Abbildung $g^*(p) : g^*(E) \to X$, wobei $g^*(E) \coloneqq X \times_B E$ das Faserprodukt von $X$ und $E$ über $B$ vermöge $g$ und $p$ ist.
\end{defn}

\begin{bem}
  Pullback ist funktoriell: $(g \circ f)^* = f^* \circ g^*$ und $\id^* = \id$.
\end{bem}

\begin{lem}
  Pullbacks von Serre-Faserungen sind Serre-Faserungen.
\end{lem}

\begin{proof}
  Sei $p : E \to B$ eine Serre-Faserung und $g : X \to B$ stetig.
  Wir müssen die Existenz des Morphismus $\tilde{H}$ im folgenden Diagramm zeigen:
  \begin{centertikzcd}[row sep=1.2cm, column sep=1.4cm]
    D^n \arrow[r, "H_0"] \arrow[d, hook, "i_0", swap] &
    g^*(E) \arrow[d, "g^*(p)"] \arrow[r, "h"] \arrow[dr, phantom, "\ulcorner", very near start] &
    E \arrow[d, "p"] \\
    D^n \times \I \arrow[r, "H"] \arrow[ur, dashed, "\tilde{H}"] &
    X \arrow[r, "g"] &
    B
  \end{centertikzcd}
  Aus der HLE von $p$ erhält wie folgt einen Morphismus $K$:
  \begin{centertikzcd}[row sep=1.2cm, column sep=1.4cm]
    D^n \arrow[r, "H_0"] \arrow[d, hook, "i_0", swap] &
    X \times_B E \arrow[r, "h"] &
    E \arrow[d, "p"] \\
    D^n \times \I \arrow[r, "H"] \arrow[urr, "K", dashed] &
    X \arrow[r, "g"] &
    B
  \end{centertikzcd}
  Nun ist $D^n \times \I$ vermöge $H$ und $K$ ein Kegel über dem Diagramm $(X \xrightarrow{g} B \xleftarrow{p} E)$.
  Die universelle Eigenschaft von $g^*(E)$ induziert einen Morphismus $\tilde{H} : D^n \times \I \to X \times_B E$ mit
  $g^*(p) \circ \tilde{H} = H$ und $h \circ \tilde{H} = K$.
  Aus der univ. Eigenschaft von $g^*(E)$ (Eindeutigkeit) folgt nun $\tilde{H} \circ i_0 = H_0$.
  % weil: $h \circ H_0 = K \circ i_0 = h \circ \tilde{H} \circ i_0$
\end{proof}

\begin{defn}
  Ein Morphismus $(g, \tilde{g}) : p' \to p$ von Serre-Faserungen $p' : E' \to B'$ und $p : E \to B$ ist ein kommutatives Quadrat der Form
  \begin{centertikzcd}
    E' \arrow[r, "\tilde{g}"] \arrow[d, "p'"] &
    E \arrow[d, "p"] \\
    B' \arrow[r, "g"] &
    B
  \end{centertikzcd}
\end{defn}

\begin{bsp}
  Pullback einer Serre-Faserung $p$ entlang einer stetigen Abbildung $g$ induziert einen Morphismus $(g, \tilde{g}) : g^*(p) \to p$ von Serre-Faserungen.
\end{bsp}

\begin{lem}
  Die langen exakten Sequenzen der Homotopiegruppen von Faserungen sind natürlich: \ES{} $(g, \tilde{g}) : p' \to p$ ein Morphismus von Serre-Faserungen $p' : E' \to B'$ und $p : E \to B$, $b_0' \in B'$, $b_0 \coloneqq g(b_0')$, $F' \coloneqq p'^{-1}(b_0')$, $F \coloneqq p^{-1}(b_0)$, $f'_0 \in F'$, $f_0 \coloneqq \tilde{g}(f'_0)$. Dann gibt es eine "`Leiter"' bestehend aus kommutativen Quadraten zwischen den Homotopiesequenzen:
  \begin{centertikzcd}[row sep=0.8cm, column sep=0.6cm]
    \ldots \arrow[r] &
    \pi_n(F', f'_0) \arrow[r, "i'_*"] \arrow[d, "(\tilde{g}|_{F'})_*"] &
    \pi_n(E', f'_0) \arrow[r, "p'_*"] \arrow[d, "\tilde{g}_*"] &
    \pi_n(B', b'_0) \arrow[r, "\partial"] \arrow[d, "g_*"] &
    \pi_{n-1}(F', f'_0) \arrow[r] \arrow[d, "(\tilde{g}|_{F'})_*"] &
    \ldots \\
    \ldots \arrow[r] &
    \pi_n(F, f_0) \arrow[r, "i_*"] &
    \pi_n(E, f_0) \arrow[r, "p_*"] &
    \pi_n(B, b_0) \arrow[r, "\partial"] &
    \pi_{n-1}(F, f_0) \arrow[r] &
    \ldots
  \end{centertikzcd}
\end{lem}

\begin{proof}
  Folgt aus der Natürlichkeit der langen exakten Homotopiesequenz von Raumpaaren.
\end{proof}

\ES{} $p : E \to B$ eine Serre-Faserung, $\gamma : \I \to B$ ein stetiger Weg.
Betrachte die lange exakte Sequenz
% TODO: Aufpunkte, Zusammenhang, H_0?
\[ \ldots \to \pi_n(F_{\gamma(0)}) \to \pi_n(\gamma^*(E)) \to \pi_n(\I) \to \pi_{n-1}(F_{\gamma(0)}) \to \ldots \]
der Homotopiegruppen von $\gamma^*(p) : \gamma^*(E) \to \I$ mit Faser
\[ F_{\gamma(t)} \coloneqq \gamma^*(p)^{-1}(t)  \subset \gamma^*(E) = \Set{(t, e) \in \I \times E}{\gamma(t) = p(e)}. \]
In dieser Sequenz sind die Gruppen $\pi_n(\I)$ trivial.
Folglich sind die Abbildungen $(i_{\gamma(t)})_* : \pi_n(F_{\gamma(t)}, *) \to \pi_n(\gamma^*(E), *)$ Isomorphismen. % TODO: Anderer Name für die Aufpunkte
In anderen Worten: $i_{\gamma(t)}$ ist eine schwache Äquivalenz.
Aus einem Korollar des Whitehead-Theorems folgt nun, dass $i_t$ auch in Homologie und Kohomologie Isomorphismen induziert (vgl. Spanier, AT, S. 406, Cor 7.6.25). % TODO: Besser referenzieren!
Wir untersuchen den Isomorphismus
\[ T_\gamma \coloneqq (i_{\gamma(1)})^* \circ ((i_{\gamma(0)})^*)^{-1} \enspace:\enspace H^*(F_{\gamma(0)}) \xrightarrow{\cong} H^*(F_{\gamma(1)}). \]

\begin{lem}
  $T_\gamma$ hängt lediglich von der Weghomotopieklasse von $\gamma$ ab, \dh{} ist $\eta$ ein zweiter Weg mit $\gamma \simeq \eta$, so gilt $T_\gamma = T_\eta$.
\end{lem}

\begin{proof}
  Sei $H : \I \times \I \to B$ eine Homotopie zw. den Wegen $\gamma$ und $\eta$, \dh{} $H_0 \coloneqq H(0, \blank) = \gamma$, $H_1 = \eta$, $H(\blank, 0) \equiv x$ und $H(\blank, 1) \equiv y$ mit $x \coloneqq \gamma(0) \!=\! \eta(0)$ und $y \coloneqq \gamma(1) \!=\! \eta(1)$.
  Für festes $s \in \I$ sei $i_s : \I \to \I \times \I, \enspace t \mapsto (s, t)$ die Inklusion als $\{ s \} \times I$.
  Betrachte das kommutative Diagramm
  \begin{centertikzcd}[row sep=1.2cm, column sep=1.4cm]
    H_s^*(E) \arrow[r, hook, "\widetilde{i_s}"] \arrow{d}{}[swap]{H_s^*(p)} \arrow[dr, phantom, "\ulcorner", very near start] &
    H^*(E) \arrow[r] \arrow[d, swap, "H^*(p)"] \arrow[dr, phantom, "\ulcorner", very near start] &
    E \arrow[d, "p"] \\
    \I \arrow[r, hook, "i_s"] \arrow[rr, bend right, "H_s"] &
    \I \!\times\! \I \arrow[r, "H"] &
    B
  \end{centertikzcd}
  Sei $t \in \I$ fest.
  Sei $F_{s,t} \coloneqq \left( H_s^*(p) \right)^{-1}(t) = \left( H^*(p) \right)^{-1}((s, t))$ und $f_0 \in F$. % TODO: Warum gibt es solch ein $f_0$?
  Das linke komm. Diagramm induziert einen Morphismus zw. den langen ex. Homotopieseq. von $H_t^*(p)$ und~$H^*(p)$:
  \begin{centertikzcd}[row sep=0.8cm, column sep=0.8cm]
    \ldots \arrow[r] &
    \pi_{n+1}(I, t) \arrow[r, "\partial"] \arrow[d, "i_{s*}"] &
    \pi_n(F_{s,t}, f_0) \arrow[r, "(i'_{s,t})*"] \arrow[d, equal] &
    \pi_n(H_s^*(E), f_0) \arrow{r}{H_s^*(p)_*}{} \arrow{d}{(\widetilde{i_s})_*}{} &
    \pi_n(\I, t) \arrow[r] \arrow{d}{i_{s*}}{} &
    \ldots \\
    \ldots \arrow[r] &
    \pi_{n+1}(I \!\times\! I, (s, t)) \arrow[r, "\partial"] &
    \pi_n(F_{s,t}, f_0) \arrow[r, "(i_{s,t})_*"] &
    \pi_n(H^*(E), f_0) \arrow[r, "H^*(p)_*"] &
    \pi_n(\I \!\times\! \I, (s, t)) \arrow[r] &
    \ldots
  \end{centertikzcd}
  In diesen Sequenzen verschwinden die Gruppen $\pi_n(I, t)$ bzw. $\pi_n(I \ntimes I, (s, t))$.
  Folglich induzieren die Abbildungen $\widetilde{i_s}$ Isomorphismen in Homotopie und in Kohomologie.
  Es gilt nun
  \begin{align*}
    T_\gamma &
    = (i'_{0,1})^* \circ ((i'_{0,0})^*)^{-1}
    = (i'_{0,1})^* \circ (\widetilde{i_0})^* \circ (\widetilde{i_0})^{-1} \circ ((i'_{0,0})^*)^{-1} \\
    & = (i_{0,1})^* \circ ((i_{0,0})^*)^{-1}
    \stackrel{(\star)}{=} (i_{1,1})^* \circ ((i_{1,0})^*)^{-1} \\
    & = (i'_{1,1})^* \circ (\widetilde{i_1})^* \circ (\widetilde{i_1})^{-1} \circ ((i'_{1,0})^*)^{-1}
    = (i'_{1,1})^* \circ ((i'_{1,0})^*)^{-1}
    = T_\eta.
  \end{align*}
  Die Gleichung $(\star)$ gilt wegen $i_{0,1} \simeq i_{1,1}$ und $i_{0,0} \simeq i_{1,0}$.
\end{proof}

Mit ganz ähnlicher Technik kann man zeigen:

\begin{lem}
  Es seien $\gamma, \eta : I \to B$ stetige Wege mit $\gamma(1) = \eta(0)$.
  Dann gilt
  \[ T_\eta \circ T_\gamma = T_{\gamma \centerdot \eta} \enspace:\enspace H^*(F_{\gamma(0)}) \xrightarrow{\cong} H^*(F_{\eta(1)}). \]
  Dabei ist die Komposition $\gamma \centerdot \eta$ von $\gamma$ und $\eta$ folgender Weg:
  \[
    \gamma \centerdot \eta : I \to B, \quad
    s \mapsto \begin{cases}
      \gamma(2s), & \text{falls } s \in [0, \tfrac{1}{2}], \\
      \eta(2s - 1), & \text{falls } s \in [\tfrac{1}{2}, 1].
    \end{cases}
  \]
\end{lem}

\begin{proof}
  Betrachte folgendes kommutatives Diagramm:
  \begin{centertikzcd}[row sep=1.2cm, column sep=1.4cm]
    \gamma^*(E) \arrow[r, hook, "\widetilde{j}"] \arrow{d}{}[swap]{\gamma^*(p)} \arrow[dr, phantom, "\ulcorner", very near start] &
    (\gamma \centerdot \eta)^*(E) \arrow[r] \arrow[d, swap, "(\gamma \centerdot \eta)^*(p)"] \arrow[dr, phantom, "\ulcorner", very near start] &
    E \arrow[d, "p"] \\
    \I \arrow[r, hook, "j"] \arrow[rr, bend right, "\gamma"] &
    \I \arrow[r, "\gamma \centerdot \eta"] &
    B
  \end{centertikzcd}
  Dabei ist $j : I \to I$ die Abbildung $s \mapsto \sfrac{s}{2}$.
  Analog zum letzten Lemma sieht man anhand des Leiterdiagramms der langen exakten Sequenzen der Faserungen $\gamma^*(p)$ und $(\gamma \centerdot \eta)^*(p)$, dass $\widetilde{j}$ einen Isomorphismus in Homotopie und Kohomologie induziert.
  Es gibt ein ähnliches Diagramm mit $\eta$ statt $\gamma$ und $k : I \to I, \enspace s \mapsto \sfrac{(1+s)}{2}$ statt $j$.
  Es induziert auch $\widetilde{k}$ einen Isomorphismus in Kohomologie.
  Es gilt nun
  \begin{align*}
    T_\eta \circ T_\gamma
    & = (i_{\eta(1)})^* \circ ((i_{\eta(0)})^*)^{-1} \circ (i_{\gamma(1)})^* \circ ((i_{\gamma(0)})^*)^{-1} \\
    & = (i_{\eta(1)})^* \circ \tilde{k}^* \circ (\tilde{k}^*)^{-1} \circ ((i_{\eta(0)})^*)^{-1} \circ (i_{\gamma(1)})^* \circ \tilde{j}^* \circ (\tilde{j}^*)^{-1} \circ ((i_{\gamma(0)})^*)^{-1} \\
    & = (\tilde{k} \circ i_{\eta(1)})^* \circ ((\tilde{k} \circ i_{\eta(0)})^*)^{-1} \circ (\tilde{j} \circ i_{\gamma(1)})^* \circ ((\tilde{j} \circ i_{\gamma(0)})^*)^{-1} \\
    & = (i_{\gamma \centerdot \eta(1)})^* \circ ((i_{\gamma \centerdot \eta(\sfrac{1}{2})})^*)^{-1} \circ (i_{\gamma \centerdot \eta(\sfrac{1}{2})})^* \circ ((i_{\gamma \centerdot \eta(0)})^*)^{-1} \\
    & = (i_{\gamma \centerdot \eta(1)})^* \circ ((i_{\gamma \centerdot \eta(0)})^*)^{-1} = T_{\gamma \centerdot \eta}. \qedhere
  \end{align*}
\end{proof}

\section{Wichtige Faserungen}

Das einfachste Beispiel einer Hurewicz-Faserung ist die \emph{Produktfaserung}, die Projektion $p_0 : X \times Y \to X$ eines Produktraumes auf einen Faktor mit Faser $Y$.

% TODO: wohin?
\begin{defn}
  Der \emph{Pfadraum} eines punktierten topologischer Raum $(X, x_0)$ ist
  \[ PX \coloneqq \Set{\gamma \in X^\I}{\gamma(0) = x_0} \subset X^I \]
  mit der Unterraumtopologie des Raumes $X^I$, welcher die Kompakt-Offen-Topologie besitzt.
  Der Basispunkt von $PX$ ist der konstante Weg $y_0 : \I \to X, \enspace t \mapsto x_0$.
\end{defn}

\begin{bem}
  Der Raum $PX$ ist zusammenziehbar: Die Abbildung
  \[
    H : \I \times PX \to PX, \quad
    (t, \gamma) \mapsto \gamma(t \cdot \blank)
  \]
  ist eine Homotopie zwischen der konstanten Abbildung mit Wert $\gamma_0$ und $\id_{PX}$.
\end{bem}

\begin{lem}
  Die Abbildung
  $
    p : PX \to X, \enspace
    \gamma \mapsto \gamma(1)
  $
  ist eine Hurewicz-Faserung.
\end{lem}

\begin{proof}
  Es sei ein topologischer Raum $A$ und stetige Abbildungen $H_0 : A \to PX$, $H : \I \times A \to X$ mit $H \circ i_0 = p \circ H_0$ gegeben.
  %Definiere $\tilde{H}$ durch
  Dann ist eine Homotopieliftung gegeben durch
  \begin{align*}
    \tilde{H} : \I \times A & \to PX, \\
    (s, a) & \mapsto \gamma_{s, a}, \quad
    \gamma_{s, a}(t) \coloneqq \begin{cases}
      H_0(a)(t \cdot (1+s)) & \text{falls $t \cdot (1+s) \leq 1$,} \\
      H(t \cdot (1+s) - 1, a) & \text{falls $t \cdot (1+s) \geq 1$.}
      \qedhere
    \end{cases}
  \end{align*}
\end{proof}

Die Faserung $p : PX \to X$ wird \emph{Pfadfaserung} genannt.

\begin{bem}
  Die Faser von $p$ über $x_0$ ist
  \[ \Omega X \coloneqq \Set{\gamma \in X^\I}{\gamma(0) = \gamma(1) = x_0}. \]
  Der Raum $\Omega X$ heißt \emph{Schleifenraum} von $X$.
\end{bem}

% XXX: $\Sigma X$ definieren?

\begin{bem}
  Es seien $(X, x_0)$ und $(Y, y_0)$ punktierte Räume.
  Es gibt eine in $X$ und $Y$ natürliche Bijektion
  % XXX: Einschränkung auf Unterkategorie der topologischen Räume nötig?
  \begin{alignat*}{4}
    \Hom((\Sigma X, x_0), (Y, y_0)) & \enspace\cong\enspace && \Hom((X, x_0), (\Omega Y, \gamma_0)), \\
    f & \enspace\mapsto\enspace && (x \mapsto t \mapsto f([(x, t)])), \\
    ([(x, t)] \mapsto g(x)(t)) & \enspace\mapsfrom\enspace && g.
  \end{alignat*}
\end{bem}

\begin{lem}\label{convert-to-fibration}
  Man kann jede stetige Abbildung $f : X \to Y$ schreiben als Komposition
  \[ X \xrightarrow{i} E_f \xrightarrow{p} Y \]
  einer Homotopieäquivalenz $i$ und einer Hurewicz-Faserung $p$.
  Genauer gilt
  \begin{align*}
    E_f & \coloneqq \Set{(x, \gamma) \in X \times X^\I}{f(x) = \gamma(0)} \subset X \times Y^\I, \\
    i(x) & \coloneqq (x, t \mapsto f(x)), \\
    p(x, \gamma) & \coloneqq \gamma(1).
  \end{align*}
\end{lem}

% Notation wie in https://en.wikipedia.org/wiki/Homotopy_fiber
\begin{proof}
  Offensichtlich sind $i$ und $p$ stetig und es gilt $p \circ i = f$.
  Das Homotopie-Inverse von $i$ ist $j : E_f \to X, \enspace (x, \gamma) \mapsto x$.
  Es gilt $j \circ i = \id_X$ und eine Homotopie zwischen $i \circ j$ und $id_{E_f}$ ist gegeben durch
  \[
    H : \I \times E_f \to E_f, \quad
    (s, (x, \gamma)) \mapsto (x, \gamma(s \cdot \blank)).
  \]
  Es bleibt zu zeigen, dass $p$ eine Faserung ist.
  Es sei dazu ein topologischer Raum $A$ und Abbildungen $H_0 : A \to E_f$ und $H : \I \times A \to Y$ mit $H \circ i_0 = p \circ H_0$ gegeben.
  Dann ist folgende Abbildung eine Homotopieliftung:
  \begin{align*}
    \tilde{H} : \I \times A & \to E_f, \\
    (s, a) & \mapsto \gamma_{s, a}, \quad
    \gamma_{s, a}(t) \coloneqq \begin{cases}
      H_0(a)(t \cdot (1+s)) & \text{falls $t \cdot (1+s) \leq 1$,} \\
      H(t \cdot (1+s) - 1, a) & \text{falls $t \cdot (1+s) \geq 1$.}
      \qedhere
    \end{cases}
  \end{align*}
\end{proof}
% TODO: Die Homotopieliftung ist zweimal genau gleich als Abbildung. Kann man die Pfadfaserung als Spezialfall vom Lemma sehen und sich das einmal sparen?

\section{Lokale Koeffizienten}

% hier nach http://isites.harvard.edu/fs/docs/icb.topic880873.files/local_systems.pdf

\begin{defn}
  Ein \emph{lokales Koeffizientensystem} $\underline{A}$ auf einem topologischen Raum $B$ besteht aus abelschen Gruppen $(A_b)_{b \in B}$ und Isomorphismen $T_\gamma : A_{\gamma(0)} \xrightarrow{\cong} A_{\gamma(1)}$ für jeden stetigen Weg $\gamma : I \to B$, sodass gilt:
  \begin{itemize}
    \item Sind zwei Wege $\gamma, \eta : I \to B$ homotop modulo Endpunkte, so gilt $T_\gamma = T_\eta$.
    \item Für komponierbare Wege $\gamma, \eta : I \to B$ gilt $T_{\gamma \centerdot \eta} = T_\eta \circ T_\gamma$.
    \item Für den konstanten Weg $\gamma \equiv b$ gilt $T_\gamma = \id_{A_b}$.
  \end{itemize}
\end{defn}

\begin{bem}
  % TODO: sagt man "lokales Koeffizientensystem" oder nur "lokales System"?
  Man kann ein lokales Koeffizientensystem auf $B$ auch als Funktor von dem Fundamentalgruppoid von $B$ in die Kategorie der abelschen Gruppen auffassen.
\end{bem}

\begin{bsp}
  Im letzten Abschnitt wurde gezeigt: Bei einer Serre-Faserung $p : E \to B$ bilden die $q$-ten Kohomologiegruppen $A_b \coloneqq H^q(p^{-1}(b))$ der Fasern ein lokales Koeffizientensystem.
  Mit dem universellen Koeffizententheorem sieht man, dass gleiches auch für die Homologiegruppen $A_b \coloneqq H^q(p^{-1}(b); G)$ mit Koeffizienten in einer abelschen Gruppe $G$ gilt.
  Wir bezeichnen dieses Koeffizientensystem im Folgenden mit $\LH^q(F_p; G)$. % XXX: Das ist ne doofe Notation. Ändern?
\end{bsp}

\begin{bsp}
  Für jede abelsche Gruppe $G$ gibt es das konstante Koeffizientensystem $\underline{G}$ mit $G_b \coloneqq G$ für alle $b \in B$ und $T_\gamma = \id_G$ für alle $\gamma : I \to B$.
\end{bsp}

Sei im Folgenden $\Delta_n(B)$ die Menge der $n$-Simplizes in $B$, also die Menge der stetigen Abbildungen $\Delta^n \to B$ mit $\Delta^n \coloneqq \spann \{ e_0, \ldots, e_n \} \subset \R^{n+1}$, und
\[
  d_n : \Delta_n(B) \to \Delta_{n-1}(B) \quad
  \sigma \mapsto \sigma_{\angles{e_0, \ldots, \hat{e_i}, \ldots, e_n}} \qquad (0 \leq i \leq n),
\]
die Abbildung auf die $i$-Seite. % TODO: kann man das "Projektion" nennen?
Für einen $n$-Simplex $\sigma$ bezeichne $\sigma_i \coloneqq \sigma_{\angles{e_i}} \in \Delta_0(B) = B$ die $i$-te Ecke und $\sigma_{ij} \coloneqq \sigma_{\angles{e_i, e_j}} \in \Delta_1(B)$ den Weg von von $\sigma_i$ nach $\sigma_j$ entlang der $ij$-Kante von $\sigma$ ($0 \leq i \leq j \leq n$).

\begin{defn}
  Sei $B$ ein topologischer Raum, $\underline{A}$ ein lokales Koeffizientensystem auf $B$.
  Der \emph{Kokomplex der singulären Koketten auf $B$ mit Koeffizienten in $\underline{A}$} ist folgendermaßen definiert:
  \[
    C^n(B; \underline{A}) \coloneqq \prod_{\mathclap{\sigma \in \Delta_n(B)}} A_{\sigma_0}, \quad
    \delta^n\left((a_\tau)_{\tau \in \Delta_n(B)}\right)_{\sigma \in \Delta_{n+1}(B)} \coloneqq T_{\sigma_{01}}^{-1}(a_{d_0(\sigma)}) + \sum_{i=1}^{n+1} (-1)^i a_{d_i(\sigma)}.
  \]
  Man überprüft leicht, dass $\delta^{n+1} \circ \delta^n = 0$ gilt.
  Die Kohomologie $H^*(B; \underline{A}) \coloneqq H^*(C^*(B; \underline{A}))$ dieses Kettenkomplexes heißt \emph{singuläre Kohomologie von $B$ mit Koeffizienten in $\underline{A}$}.
\end{defn}

\begin{beob}
  Für das konstante Koeffizientensystem $\underline{G}$ gilt $H^*(B; \underline{G}) \cong H^*(B; G)$.
  Gewöhnliche Kohomologie mit Koeffizienten ist also ein Spezialfall von Kohomologie mit Koeffizienten in einem lokalen System.
\end{beob}

\begin{defn}
  Es sei $\underline{R}$ ein lokales Koeffizientensystem, in dem die Gruppen $R_b$ sogar Ringe und die Abbildungen $T_\gamma$ Ringisomorphismen sind. Dann definiert
  \begin{align*}
    & \cup : H^m(B; \underline{R}) \times H^n(B; \underline{R}) \to H^{m+n}(B; \underline{R}), \\
    & ([(a_\sigma)_{\sigma \in \Delta_m(B)}] \cup [(b_\sigma)_{\sigma \in \Delta_n(B)}])_{\tau \in \Delta_{m+n}(B)} \coloneqq a_{\sigma_{\angles{e_1, \ldots, e_m}}} \cdot T_{\sigma_{0m}}^{-1}(b_{\sigma_{e_m, \ldots, e_{m+n}}})
  \end{align*}
  ein Produkt, das sogenannte \emph{Cup-Produkt}.
\end{defn}

\section{Spektralsequenzen}

Ein häufig verwendetes Hilfsmittel in der algebraischen Topologie sind lange exakte Sequenzen.
Sie liefern einen Zusammenhang zwischen Homologie- oder Homotopiegruppen von verschiedenen Räumen.
Wenn man genügend viele dieser Gruppen kennt, so kann man oft rein algebraisch die anderen Gruppen erschließen.
Ein Beispiel ist die lange exakte Homotopiesequenz einer Faserung.
Nun kann man sich fragen, ob es bei Faserungen auch einen Zusammenhang zwischen den Homologie- und Kohomologiegruppen von Basisraum, Totalraum und Faser gibt.
Im Jahr 1951 hat Jean-Pierre Serre in seiner Dissertation \cite{serre:thesis} gezeigt, dass es einen solchen tatsächlich gibt.
Dieser hat jedoch nicht die Form einer langen exakten Sequenz, sondern ist kodiert in einem komplizierten algebraischen Objekt, einer Spektralsequenz.

\ES{} $A$ im Folgenden ein kommutativer Ring mit Eins.

\begin{defn}
  Eine (kohomologische) \emph{Spektralsequenz} besteht aus
  \begin{itemize}
    \item $A$-Moduln $E_r^{p,q}$ für alle $p, q \in \Z$ und $r \geq 1$,
    \item $A$-Modul-Homomorphismen $d_r^{p,q} : E_r^{p,q} \to E_r^{p+r,q-r+1}$ mit $d_r^{p+r,q-r+1} \circ d_r^{p,q} = 0$
    \item und Isomorphismen $\alpha_r^{p,q} : H^{p,q}(E_r) \!\coloneqq\! \ker(d_r^{p,q}) / \im(d_r^{p-r,q+r-1}) \xrightarrow{\cong} E_{r+1}^{p,q}$.
  \end{itemize}
\end{defn}

\begin{bem}
  \begin{itemize}
    \item Die Homomorphismen $d^r_{p,q}$ heißen \emph{Differentiale}.
    \item Die Gesamtheit der Moduln $E_r^{p,q}$ und Differentiale $d_r^{pq}$ mit $r \in \N$ fest heißt \emph{$r$-te Seite} $E_r$.
    \item Die Isomorphismen $\alpha_r^{p,q}$ werden in der Notation  unterdrückt und so gerechnet, als wäre $E^{p,q}_{r+1}$ gleich $H^{p,q}(E_r)$.
    Wenn man eine Seite einer Spektralsequenz kennt, kann man also die Einträge auf der nächsten Seite berechnen, die Differentiale jedoch im Allgemeinen nicht.
  \end{itemize}
  Man stellt Seiten für gewöhnlich in einem 2-dimensionalen Raster dar:
  \begin{center}
    \begin{tikzpicture}[x=12pt,y=12pt]
      \begin{scope}[shift={(0,0)}]
        \foreach \x in {-1,...,4}{
          \foreach \y in {-2,-1,...,3}{
            \zeroDot{\x}{\y}
          }
        }
        \foreach \x in {-1,...,5}{
          \foreach \y in {-2,-1,...,3}{
            \draw[->,gray] (\x-0.8,\y) -- (\x-0.2,\y);
          }
        }
        \draw[->] (-0.35,-0.35) -- (5.35,-0.35) node[below] {p};
        \draw[->] (-0.35,-0.35) -- (-0.35,4) node[left] {q};
        \node at (5,4.5) {$E_1$};
      \end{scope}
      \begin{scope}[shift={(10,0)}]
        \foreach \x in {-1,...,4}{
          \foreach \y in {-2,-1,...,3}{
            \zeroDot{\x}{\y}
          }
        }
        \foreach \x in {0,...,5}{
          \foreach \y in {-2,-1,...,3}{
            \draw[->,gray] (\x-1.8,\y+0.9) -- (\x-0.2,\y+0.1);
          }
        }
        \draw[->] (-0.35,-0.35) -- (5.35,-0.35) node[below] {p};
        \draw[->] (-0.35,-0.35) -- (-0.35,4) node[left] {q};
        \node at (5,4.5) {$E_2$};
      \end{scope}
      \begin{scope}[shift={(20,0)}]
        \foreach \x in {-1,...,4}{
          \foreach \y in {-2,-1,...,3}{
            \node[draw,circle,inner sep=0.5pt,fill] at (\x,\y) {};
          }
        }
        \foreach \x in {1,...,5}{
          \foreach \y in {-2,-1,...,2}{
            \draw[->,gray] (\x-2.8,\y+1.9) -- (\x-0.2,\y+0.1);
          }
        }
        \draw[->] (-0.35,-0.35) -- (5.35,-0.35) node[below] {p};
        \draw[->] (-0.35,-0.35) -- (-0.35,4) node[left] {q};
        \node at (5,4.5) {$E_3$};
      \end{scope}
    \end{tikzpicture}
  \end{center}
\end{bem}

\begin{defn}
  Eine Spektralsequenz \emph{konvergiert}, falls für alle $p, q \in \Z$ ein $R \in \N$ existiert, sodass für alle $r \geq R$ die Differentiale von und nach $E_r^{p,q}$ verschwinden und damit $E^\infty_{p,q} \coloneqq E_R^{p,q} \cong E_{R+1}^{p,q} \cong \ldots$ \\[2pt]
  Der Grenzwert der SS ist die Unendlich-Seite $E_\infty \coloneqq \{ E_\infty^{p,q} \}_{p,q}$.
\end{defn}

\begin{bem}
  Viele Spektralsequenzen sind im ersten Quadranten konzentriert, \dh{} $E_r^{p,q}$ ist nur für $p, q \geq 0$ ungleich Null.
  Solche Spektralsequenzen konvergieren immer, denn für alle $p, q \in \Z$ führen für $r \geq \max(p+1, q+2)$ alle Differentiale von $E_r^{p,q}$ aus dem ersten Quadranten heraus und alle dort eintreffenden Differentiale kommen von außerhalb des ersten Quadranten und sind daher Null.
\end{bem}

\begin{defn}
  Eine \emph{Filtrierung} eines $A$-Moduls $M$ ist eine absteigende Folge
  \[ M \supseteq \ldots \supseteq F^{p-1} M \supseteq F^p M \supseteq F^{p+1} M \supseteq \ldots \]
  von Untermoduln von $M$, $p \!\in\! \Z$.
  Eine Filtrierung heißt
  \begin{itemize}
    \item \emph{ausschöpfend}, falls $M = \bigcup_p F^p$,
    \item \emph{Hausdorffsch}, wenn $0 = \bigcap_p F^p M$ und
    \item \emph{regulär}, wenn sie ausschöpfend und Hausdorffsch ist.
  \end{itemize}
\end{defn}

\begin{defn}
  Eine Spektralsequenz $E$ \emph{konvergiert gegen} einen graduierten $A$-Modul $M = \oplus_{n \in \Z} M^n$ (notiert $E_r^{p,q} \Rightarrow M^{p+q}$), falls $E$ überhaupt konvergiert und reguläre Filtrierungen
  \[ M^n \supseteq \ldots \supseteq F^{p-1} M^n \supseteq F^p M^n \supseteq F^{p+1} M^n \supseteq \ldots \]
  existieren, sodass $E_\infty^{pq} \cong F^p M^{p+q} / F^{p+1} M^{p+q}$ für alle $p, q \in \Z$.
\end{defn}

In der Anwendung von Spektralsequenzen kennt man oft die Einträge der $E_1$ oder $E_2$-Seite.
Man versucht dann, die Einträge der $E_\infty$-Seite zu bestimmen, indem man sukzessive die Seiten $E_r$ ausrechnet.
Das Problem dabei ist, dass die dazu benötigten Differentiale auf diesen Seiten im Allgemeinen nicht bekannt sind.
In der Praxis hofft man darauf, dass Differentiale allein deswegen verschwinden, da ihre Quell- bzw. Zielgruppe die Nullgruppe ist.
Man kann auch partielle Information über die $E_\infty$-Seite verwenden, um Aussagen über Differentiale zu treffen, z.\,B. dass ein Differential ein Isomorphismus ist.

Angenommen, man hat auf diesem Weg die Seite $E_\infty$ einer Spektralsequenz, welche gegen $M$ konvergiert, berechnet.
Dann kennt man immer noch nicht die Moduln $M^n$, sondern lediglich die Quotienten in einer Filtrierung von $M^n$.
Im Fall, dass $A$ ein Körper ist, ist $M^n$ isomorph zur direkten Summe dieser Quotienten.
Im Allgemeinen ist das aber nicht der Fall, wie das Beispiel der Filtrierung $\Z \supseteq 2 \Z \supseteq 0$ der abelschen Gruppe $\Z$ zeigt.

Anders herum kann man Spektralsequenzen auch verwenden, um von $E_\infty$ (oder den Moduln $M^n$) auf $E_2$ bzw. $E_1$ zu schließen.
Dies werden wir tun, um die Homologie von Eilenberg-MacLane-Räumen mit Koeffizienten in $\Q$ zu bestimmen.

\section{Die Spektralsequenz eines filtrierten Komplexes}

% Hier nach http://people.mpim-bonn.mpg.de/matschke/SpectralSequences.pdf

\begin{defn}
  Eine \emph{Filtrierung eines Kokettenkomplexes} $C^\bullet$ ist eine absteigende Folge
  \[ C^\bullet \supseteq \ldots \supseteq F^{p-1} C^\bullet \supseteq F^p C^\bullet \supseteq F^{p+1} C^\bullet \supseteq \ldots \]
 von Unterkomplexen.
\end{defn}

% TODO: Morphismus von filtrierten Komplexen?

\begin{lem}
  \ES{} $C^\bullet$ ein filtrierter Kokettenkomplex.
  Es gibt eine Spektralsequenz mit
  \[ E_1^{pq} = H^{p+q}(F^p C^\bullet / F^{p+1} C^\bullet). \]
  Angenommen, die Filtrierung ist
  \begin{enumerate}[label=\alph*)]
    \item \emph{gradweise nach unten beschränkt}, \dh{} für alle $q \in \Z$ gibt es ein $p \in \Z$ mit $F^p C^q = 0$,
    \item \emph{ausschöpfend}, \dh{} für alle $q \in \Z$ ist $\cup_p F^p C^q = C^q$ und %(aus der ersten Bedingung folgt schon $\cap_p F^p C^q = 0$).
    \item für alle $q \in \Z$ gibt es ein $P \in \Z$, sodass für alle $p \leq P$ gilt: Die Inklusion $F^p C^\bullet \hookrightarrow C^\bullet$ induziert einen Isomorphismus $H^q (F^p C^\bullet) \cong H^q(C^\bullet)$ in Kohomologie.
  \end{enumerate}
  Dann konvergiert die Spektralsequenz gegen $H^*(C^\bullet)$.
\end{lem}

Wir führen zunächst etwas neue Notation ein.
Diese hilft, den Beweis verständlicher zu formulieren.
Wir fassen im Folgenden den Kettenkomplex als ein einziges Modul $C \coloneqq \oplus_{n \in \Z} C^n$ anstatt als Folge von Moduln auf.
Dieses Modul ist filtriert durch die Untermodule $F^p \coloneqq \oplus_{n \in \Z} F^p C^n$.
Wir setzen $F^{-\infty} \coloneqq C$ und $F^{\infty} \coloneqq 0$.
Die Korandabbildung fassen wir als Homomorphismen $d : C \to C$ mit $d \circ d = 0$ auf, der die Filtrierung von $C$ respektiert.

Wir sind interessiert an der Kohomologie von $C^\bullet$, also an $H^*(C) \coloneqq \ker(d) / \im(d)$ und an der Kohomologie von $F^p / F^{p+1}$, also $H^*(F^p / F^{p+1}) \cong (d|_{F^p})^{-1}(F^{p+1}) / d(F^p)$.
Wir geben nun eine Verallgemeinerung der Definition der Kohomologie von $C^\bullet$ und der Kohomologie des Quotientenkomplexes $F^p / F^q$: Statt Zykeln (\dh{} Elementen $c \in C$ mit $d(c) = 0$) betrachten wir \emph{$z$-Zykel}, das sind Elemente $c \in C$ mit $d(c) \in F^z$. Wir teilen diese durch die Menge $d(F^b)$ der \emph{$b$-Ränder} anstatt durch die Menge $d(C)$ der Ränder. Wir setzen

\[ S[z, q, p, b] \coloneqq \frac{F^p \cap d^{-1}(F^z)}{(F^p \cap d^{-1}(F^z)) \cap (F^q + d(F ^b))}. \]

% TODO: Bemerkung zur grafischen Notation von Matschke

Wir haben als Spezialfälle
\[
  S[p,q,p,q] \cong F^p / F^q
  \quad \text{und} \quad
  S[q,q,p,p] \cong H^*(F^p / F^q).
\]

\begin{lem}\label{differential-homomorphism}
  \ES{} $z_1 \geq q_1 \geq p_1 = z_2 \geq b_1 = q_2 \geq p_2 \geq b_2$.
  Dann ist folgende Abbildung ein wohldefinierter Homomorphismus:
  \[
    d^* : S[z_2, q_2, p_2, b_2] \to S[z_1, q_1, p_1, b_1], \quad
    [c] \mapsto [d(c)].
  \]
\end{lem}

\begin{proof}
  Falls $[c] = 0$ in $S[z_2, q_2, p_2, b_2]$, so existieren $x \in F^{q_2}$ und $y \in F^{b_2}$ mit $c = x + d(y)$. Somit gilt $d^*[c] = [dc] = [d(x) + d^2(y)] = [d(x)] = 0$ in $S[z_1, q_1, p_1, b_1]$, da $F^{b_1} = F^{q_2}$.
\end{proof}

\begin{lem}\label{differential-isomorphism}
  Es seien Filtrierungsindizes wie folgt gegeben:
  \begin{centertikzcd}[row sep=0.2cm, column sep=0.5cm]
    &&&& z_3 \arrow[r, "\geq" description, phantom] \arrow[d, equal] &
    q_3 \arrow[r, "\geq" description, phantom] \arrow[d, equal] &
    p_3 \arrow[r, "\geq" description, phantom] &
    b_3 \\

    && z_2 \arrow[r, "\geq" description, phantom] \arrow[d, equal] &
    q_2 \arrow[r, "\geq" description, phantom] \arrow[d, equal] &
    p_2 \arrow[r, "\geq" description, phantom] &
    b_2 \\

    z_1 \arrow[r, "\geq" description, phantom] &
    q_1 \arrow[r, "\geq" description, phantom] &
    p_1 \arrow[r, "\geq" description, phantom] &
    b_1
  \end{centertikzcd}
  Dann ist
  \[
    \alpha : S[q_1, q_2, p_2, p_3] \to
    \frac{\ker(d^* : S[z_2, q_2, p_2, b_2] \to S[z_1, q_1, p_1, b_1])}{\im(d^* : S[z_3, q_3, p_3, b_3] \to S[z_2, q_2, p_2, b_2])}, \quad
    [c] \mapsto [c]
  \]
  ein wohldefinierter Isomorphismus.
\end{lem}

\begin{proof}
  Sei $A$ der Quotient auf der rechten Seite. \\[2pt]
  \emph{Wohldefiniertheit}: Sei $[c] = 0$ in $S[q_1, q_2, p_2, p_3]$, \dh{} es gibt $e \in F^{q_2} = F^{b_1}$ und $f \in F^{p_1}$ mit $c = e + d(f)$. Dann ist $d^*[c] = [d(c)] = [d(e)] = 0$ in $S[z_1, q_1, p_1, b_1]$, also $c \in \ker(d^* : S[z_2, q_2, p_2, b_2] \to S[z_1, q_1, p_1, b_1])$.
  Nun ist $f \in d^{-1}(F^{z_3})$, da $d(f) = c - e \in F^{p_2} = F^{z_3}$.
  Es gilt $[c] = [e + d(f)] = [d(f)] = d^*[f] = 0$ in $A$. \\[2pt]
  \emph{Injektivität}: Sei $c \in F^{p_2} \cap d^{-1}(F^{q_1})$ mit $[c] = 0$ in $A$.
  Das heißt, es gibt $e \in F^{q_2}$, $f \in F^{b_2}$ und $g \in F^{p_3} \cap d^{-1}(F^{z_3})$ mit $c = e + d(f) + d(g)$.
  Dann ist $[c] = [e + d(f+g)] = 0$ in $S[q_1, q_2, p_2, p_3]$, da $f+g \in F^{p_3}$. \\[2pt]
  \emph{Surjektivität}: Sei $\tilde{c} \in F^{p_2} \cap d^{-1}(F^{z_2})$ mit $[\tilde{c}] \in \ker(d^* : S[z_2, q_2, p_2, b_2] \to S[z_1, q_1, p_1, b_1])$.
  Das heißt, es gibt $e \in F^{q_1}$ und $f \in F^{b_1} = F^{q_2}$ mit $d(\tilde{c}) = e + d(f)$.
  Dann ist $[\tilde{c}] = [\tilde{c} - f]$ in $S[q_1, q_2, p_2, p_3]$ mit $\tilde{c} - f \in F^{p_2} \cap d^{-1}(F^{q_1})$, da $d(\tilde{c} - f) = e \in F^{q_1}$.
\end{proof}

\begin{proof}[Beweis des Lemmas über Existenz der Spektralsequenz]
  Wir beachten jetzt wieder, dass $C$ und damit $S[z, q, p, b]$ graduiert und $d$ ein Differential vom Grad $+1$ ist.
  Es sei $S[z, q, p, b]^n$ die $n$-te Komponente.
  Setze
  \[ E_r^{pq} \coloneqq S[p\!+\!r, p\!+\!1, p, p\!-\!r\!+\!1]^{p+q}. \]
  Die Differentiale sind
  \[
    d_r^{pq}
    \enspace : \enspace
    \underbrace{S[p\!+\!r, p\!+\!1, p, p\!-\!r\!+\!1]^{p+q}}_{= E_r^{p,q}} \to \underbrace{S[p\!+\!2r, p\!+\!r\!+\!1, p\!+\!r, p\!+\!1]^{p+q+1}}_{= E_r^{p+r,q-r+1}}, \quad
    [c] \mapsto [d(c)].
  \]
  Sie sind wohldefiniert nach Lemma~\ref{differential-homomorphism}.und
  wegen Lemma~\ref{differential-isomorphism} ist
  \[
    \alpha_r^{pq} : H^{p,q}(E_r) = \ker(d^{pq}_r) / \im(d^{p-r,q+r-1}_r) \to E_{r+1}^{pq}, \quad
    [c] \mapsto [c]
  \]
  ein wohldefinierter Isomorphismus.

  \emph{Beweis der Konvergenz}:
  Es seien $p, q \in \Z$.
  Wegen Bedingung a) gibt es ein $R_1 \geq 0$, sodass $F^{p+R_1} C^{p+q+1} = 0$.
  Für $r \geq R_1$ ist damit $E^{p+r,q-r+1}_r$ als Subquotient (\dh{} Quotient eines Untermoduls) von $F^{p+R_1} C^{p+q+1}$ Null.
  Folglich verschwindet auch das Differential $d_r^{pq}$.
  Wegen Bedingung c) gibt es ein $S \in \Z$, sodass $F^s C^\bullet \hookrightarrow C^\bullet$ und somit auch $F^s C^\bullet \hookrightarrow F^{s-1} C^\bullet$ für $s \leq S$ einen Isomorphismus in $H^{p+q-1}$ und $H^{p+q}$ induziert.
  Anhand der langen exakten Sequenz zu $0 \to F^s C^\bullet \to F^{s-1} C^\bullet \to F^{s-1} C^\bullet / F^s C^\bullet \to 0$ sieht man, dass $H^{p+q-1}(F^{s-1} C^\bullet / F^s C^\bullet) = 0$.
  Somit ist $E_r^{p-r,q+r-1}$ für $r \geq R_2 \coloneqq p - s + 1$ als Submodul von $H^{p+q-1}(F^{p-r} C^\bullet / F^{p-r+1} C^\bullet)$ Null.
  Folglich verschwindet auch $d_r^{p-r,q+r-1}$.
  Mit $R \coloneqq \max(R_1, R_2)$ gilt dann $E^{pq}_R \cong E^{pq}_{R+1} \cong \ldots \cong E^{pq}_\infty$.

  Sei $H^n(C^\bullet)$ absteigend filtriert durch $F^p H^n(C^\bullet) \coloneqq \im(i^* : H^n(F^p C^\bullet) \to H^n(C^\bullet))$.
  Für $r \geq R$ ist
  \[ E^{pq}_\infty \cong E^{pq}_r = \frac{F^p C^{p+q} \cap d^{-1}(0)}{(F^p C^{p+q} \cap d^{-1}(0)) \cap (F^{p+1} C^{p+q} + d(F^{p-r+1} C^{p+q-1}))} = S[\infty,p+1,p,p-r+1]^{p+q}. \]
  Es ist daher $F^p H^{p+q}(C^\bullet) / F^{p+1} H^{p+q}(C^\bullet) \cong S[\infty, p+1, p, -\infty]^{p+q}$ ein Quotient von $E^{pq}_\infty$.
  Tatsächlich gilt $S[\infty, p+1, p, -\infty]^{p+q} \cong E^{pq}_\infty$, denn:
  Sei $c \in F^p C^{p+q} \cap d^{-1}(0)$ mit $[c] = 0$ in $S[\infty, p+1, p, -\infty]^{p+q}$.
  Dann gibt es ein $e \in F^{p+1} C^{p+q}$ und ein $f \in C^{p+q-1}$ mit $c = e + d(f)$.
  Wegen Bedingung b) gibt es ein $\tilde{p} \in \Z$ mit $f \in F^{\tilde{p}} C^{p+q+1}$.
  Wähle $r$ so, dass $r \geq R$ und $p-r+1 \leq \tilde{p}$.
  Dann ist $[c] = [e]+[d(f)] = 0$ in $E^{pq}_r \cong E^{pq}_\infty$.
\end{proof}


\section{Die Serre-Spektralsequenz}

\begin{satz}[Jean-Pierre Serre]
  Es sei $G$ eine abelsche Gruppe.
  Für jede Serre-Faserung $p : E \to B$ existiert eine Spektralsequenz mit
  \[ E_2^{p,q} = H^p(B; \LH^q(F_p; G)), \]
  welche gegen $H^*(E)$ konvergiert.
\end{satz}

\begin{proof}
  \TODO{}
\end{proof}

\section{Multiplikative Struktur der Serre-Spektralsequenz}

\begin{satz}
  Es sei $R$ ein Ring, $p : E \to B$ eine Serre-Faserung.
  Dann gibt es bilineare Abbildungen
  \[
    m_r = m_r^{p,q,s,t} \enspace:\enspace E^{p,q}_r \times E^{s,t}_r \to E^{p+s,q+t}_r, \quad
    (x, y) \mapsto m_r(x, y) =: xy
  \]
  mit folgenden Eigenschaften:
  \begin{enumerate}[label=(\roman*)]
    \item $d_r$ ist derivativ: $d_r^{p+s,q+t}(xy) = (d_r^{p,q} x) y + (-1)^{p+q} x (d_r^{s,t} y)$ für alle $x \in E^{p,q}_r$ und $y \in E^{s,t}_r$
    \item Es gilt $m_{r+1}([x], [y]) = [m_r(x, y)]$ für alle $x \in \ker(d_r^{p,q})$ und $y \in \ker(d_r^{s,t})$. \\
    Dabei ist $m_r(x, y) \in \ker(d_r^{p+s,q+t})$ wegen (i).
    \item Auf der $E_2$-Seite ist $m_2 : E_2^{p,q} \!\times\! E_2^{r,s} \!\to\! E_2^{p+s,q+t}$ das $(-1)^{qs}$-fache des Cup-Produkts
    \[
      H^p(B; \LH^q(F; R)) \times H^s(B; \LH^t(F; R)) \to H^{p+s}(B; \LH^{q+t}(F; R)),
    \]
    welches für $a = [(a_\sigma)_{\sigma \in \Delta_m(B)}] \in H^p(B; \LH^q(F; R))$ und $b = [(b_\sigma)_{\sigma \in \Delta_n(B)}] \in H^s(B; \LH^t(F; R))$ definiert ist durch
    \[
      (a \cup b)_{\sigma \in \Delta_{p+s}(B)} \coloneqq
      a_{\sigma_{\angles{e_1, \ldots, e_m}}} \cup T_{\sigma_{0m}}^{-1}(b_{\sigma_{\angles{e_m, \ldots, e_{m+n}}}}).
    \]
    \item Das Cup-Produkt auf $H^*(B; R)$ respektiert die Filtrierungen von $H^n(B; R)$ und schränkt daher ein zu Abbildungen $F_p^m \times F_s^n \to F_{p+s}^{m+n}$.
    Die induzierte Abbildung auf dem Quotienten $F^m_p/F^m_{p+1} \times F^n_s/F^n_{s+1} \to F^{m+n}_{p+s} / F^{m+n}_{p+s+1}$ entspricht dem Grenzwert $m_\infty : E_\infty^{p,m-p} \times E_\infty^{s,n-s} \to E_\infty^{p+s,m+n-p-s}$ der Multiplikationen $m_r$.
  \end{enumerate}
\end{satz}

\begin{proof}
  \TODO{}
\end{proof}

% Spanier, Algebraic Topology: SS zum filtrierten Komplex; nur triviale Wirkung; etwas genauer als Hatcher
% Hatcher: nur triviale Wirkung
% Switzer, Homology and Homotopy: nur triviale Wirkung, p. 350
% May/Ponto, More Concise AT: Allgemeine Serre-Spektralsequenz (mit lokalen Koeff'en); nur ein "sketch proof"
% McCleary, User's Guide to SS: Lokale Koeffizienten, statement: p. 139, proof: p. 167ff

\section{Die Serre-Spektralsequenz für Homologie}

Es gibt eine Version des Satzes für Serre für Homologie anstatt Kohomologie.
In Homologie wird eine andere Indizierung für Spektralsequenzen verwendet:

\begin{defn}
  Eine homologische \emph{Spektralsequenz} besteht aus
  \begin{itemize}
    \item $A$-Moduln $E^r_{p,q}$ für alle $p, q \in \Z$ und $r \geq 1$,
    \item $A$-Modul-Homomorphismen $d^r_{p,q} : E^r_{p,q} \to E^r_{p-r,q+r-1}$ mit $d^r_{p-r,q+r-1} \circ d^r_{p,q} = 0$
    \item und Isomorphismen $\alpha^r_{p,q} : H_{p,q}(E^r) \!\coloneqq\! \ker(d ^r_{p,q}) / \im(d^r_{p+r,q-r+1}) \xrightarrow{\cong} E^{r+1}_{p,q}$.
  \end{itemize}
\end{defn}

Jede homologische Spektralsequenz $E$ liefert eine kohomologische Spektralsequenz, wenn man $E_r^{p,q} \coloneqq E^r_{-p,-q}$ setzt.

\begin{defn}
  Eine homologische Spektralsequenz $E$ \emph{konvergiert gegen} einen graduierten $A$-Modul $M = \oplus_{n \in \Z} M_n$ (notiert $E^r_{p,q} \Rightarrow M_{p+q}$), falls $E$ überhaupt konvergiert und reguläre Filtrierungen
  \[ 0 \subseteq \ldots \subseteq F^{p-1} M_n \subseteq F^p M_n \subseteq F^{p+1} M_n \subseteq \ldots \]
  existieren, sodass $E^\infty_{pq} \cong F^p M_{p+q} / F^{p-1} M_{p+q}$ für alle $p, q \in \Z$.
\end{defn}

\begin{satz}[Jean-Pierre Serre]
  Es sei $G$ eine abelsche Gruppe.
  Für jede Serre-Faserung $p : E \to B$ existiert eine (homologische) Spektralsequenz mit
  \[ E^2_{p,q} = H_p(B; \LH_q(F_p; G)), \]
  welche gegen $H_*(E)$ konvergiert.
\end{satz}

% TODO: Verweis auf Beweis in irgendeinem Buch.

Dabei bilden die $q$-ten Homologiegruppen der Fasern ein lokales Koeffizientensystem mit $\LH_q(F_p; G)_b = H_q(p^{-1}(b))$.
Homologie mit einem Koeffizientensystem $\underline{A}$ ist analog zu Kohomologie mit Koeffizientensystem definiert.
Falls $B$ einfach zusammenhängend ist oder allgemeiner $\pi_1(B)$ trivial auf den Homologiegruppen der Faser wirkt, so gilt $H_p(B; \LH_q(F_p; G)) \cong H_p(B; H_q(F_p; G))$.

\section{Töten von Homotopiegruppen und Eilenberg-MacLane-Räume}

\begin{defn}
  Sei $G$ eine Gruppe, $n \geq 1$.
  Ein \emph{Eilenberg-MacLane-Raum} vom Typ $K(G, n)$ ist ein punktierter, zusammenhängender topologischer Raum $(X, x_0)$ mit
  \[
    \pi_q(X, x_0) = \begin{cases}
      G & \text{falls $q = n$,} \\
      0 & \text{falls $q \neq n$.}
    \end{cases}
  \]
\end{defn}

\begin{lem}
  Sei $G$ eine abelsche Gruppe, $n \geq 2$.
  Dann existiert ein CW-Komplex $(X, x_0)$ vom Typ $K(G, n)$.
\end{lem}

\begin{proof}
  \TODO{}
\end{proof}

% TODO: Eindeutigkeit von K(G, n)'s?

\begin{bem}
  Sei $(X, x_0)$ ein $K(G, n)$. Dann ist $\Omega X$ ein $K(G, n-1)$, denn

  \begin{align*}
    \pi_q(\Omega X, \gamma_0) & \cong \Hom((S^q, *), (\Omega X, \gamma_0))
    \cong \Hom((\Sigma S^q, *), (X, x_0))
    \cong \pi_{q+1}(X, x_0) \\
    & \cong \begin{cases}
      G & \text{falls $q+1 = n \iff q = n-1$} \\
      0 & \text{falls $q+1 \neq n$.}
    \end{cases}
  \end{align*}
  % XXX: Verträglichkeit mit Gruppenstruktur???
\end{bem}

\section{Das Hurewicz-Mod-$\SC$-Theorem}

Sei $(X, x_0)$ ein punktierter topologischer Raum.
Für $n \geq 1$ liefert der Hurewicz-Homomorphismus $h_n : \pi_n(X, x_0) \to H_n(X; \Z)$ einen Zusammenhang zwischen der $n$-ten Homotopiegruppe und der $n$-ten Homologiegruppe von $X$.
Er ist definiert durch $h_n([f]) \coloneqq H_n(f)(\alpha)$ für einen fest gewählten Erzeuger $\alpha \in H_n(S^n; \Z)$.

\begin{satz}[Hurewicz]
  Sei $(X, x_0)$ ein $(n{-}1)$-zusammenhängender topologischer Raum, \dh{} $\pi_i(X, x_0) = 0$ für $i < n$.
  Dann ist $h_i : \pi_i(X, x_0) \to H_i(X; \Z)$ ein Isomorphismus für $0 < i \leq n$.
  Insbesondere gilt $H_i(X; \Z) = 0$ für $0 < i < n$.
\end{satz}

Ein Beweis dieses Satzes wird in~\cite[S. 366ff]{hatcher:at} geführt.

\begin{defn}
  Eine Klasse $\SC$ von abelschen Gruppen heißt \emph{Serre-Klasse}, falls
  \begin{enumerate}
    \item Für jede kurze exakte Sequenz $0 \to A \to B \to C \to 0$ von abelschen Gruppen gilt: $B \in \SC \iff A, B \in \SC$.
    \item Für $A, B \in \SC$, sind auch $A \otimes B \in \SC$ und $\Tor(A, B) \in \SC$.
  \end{enumerate}
\end{defn}

\begin{bem}
  Aus der ersten Eigenschaft folgt, dass Bilder, Untergruppen und Quotienten einer Gruppe aus $\SC$ wieder in $\SC$ sind.
  Genauer gilt für eine ab. Gruppe $B$ und eine Untergruppe $A < B$: $B \in \SC \iff A, B/A \in \SC$.
  Durch Induktion kann man zeigen, dass für eine Gruppe $A$ mit endlicher Filtrierung
  $A = F^0 A \supseteq F^1 A \supseteq \ldots \supseteq F^k A = 0$
  gilt: $A \in \SC \iff F^0 A / F^1 A, \ldots, F^{k-1} A / F^k A \in \SC$.
  Außerdem ist die direkte Summe zweier Gruppen aus $\SC$ wieder in $\SC$.
\end{bem}

\begin{defn}
  Es sei $\SC$ eine Serre-Klasse.
  Ein Morphismus $f : A \to B$ zwischen abelschen Gruppen heißt \emph{Isomorphimus modulo $\SC$}, falls $\ker(f), \coker(f) \in \SC$. \\
\end{defn}

\begin{bem}
  Dies ist äquivalent zur Existenz einer exakten Sequenz $K \to A \xrightarrow{f} B \to C$ mit $K, C \in \SC$.
  Eine Gruppe, welche modulo-$\SC$-isomorph zu einer Gruppe aus $\SC$ ist, ist selbst in $\SC$.
\end{bem}

\begin{bspe}
  Man kann leicht zeigen, dass folgende Klassen die Definition erfüllen:
  % XXX: Nummer der Definition
  \begin{itemize}
    \item $\FG \coloneqq \{\, \text{endlich erzeugte Gruppen} \,\}$
    \item $\F \coloneqq \{\, \text{endliche Gruppen} \,\}$
    % TODO: ist das tatsächlich leicht zu zeigen?
    % XXX: weitere Beispiele?
  \end{itemize}
\end{bspe}


% Hatcher, 1.8
\begin{satz}["`Hurewicz-mod-$\SC$-Theorem"']\label{hurewicz-mod-c}
  Es sei $\SC \in \{ \FG, \F \}$ und
  $(X, x_0)$ ein einfach zusammenhängender topologischer Raum.
  Angenommen, $\pi_i(X, x_0) \in \SC$ für $0 < i < n$.
  Dann ist $h_i : \pi_i(X, x_0) \to H_i(X; \Z)$ ein Isomorphismus modulo $\SC$ für $i \leq n$.
  Insbesondere gilt $H_i(X; \Z) \in \SC$ für $0 < i < n$.
\end{satz}

%Mit $\SC \coloneqq \{ 0 \}$ folgt das gewöhnliche Hurewicz-Theorem.

% Hatcher, 1.7 (etwas allgemeiner)
\begin{kor}\label{homotopy-in-c-iff-homology-in-c}
  Es sei $\SC \in \{ \FG, \F \}$ und
  $(X, x_0)$ ein einfach zusammenhängender topologischer Raum.
  Dann gilt für alle $N \in \N \cup \{ \infty \}$:
  \[
    \fa{0 \leq n < N} \pi_n(X, x_0) \in \SC
    \quad \iff \quad
    \fa{1 \leq n < N} H_n(X; \Z) \in \SC.
  \]
\end{kor}

\begin{proof}
  Die Aussage folgt aus Satz \ref{hurewicz-mod-c} durch Induktion über $N$.
\end{proof}

Aus dem Korollar folgt, dass die Homotopiegruppen der Sphären alle endlich erzeugt sind, da die Homologiegruppen der Sphären endlich erzeugt sind.

\begin{defn}
  Es sei $\SC$ eine Serre-Klasse.
  Wir nennen einen topologischen Raum~$X$ \emph{$\SC$-azyklisch}, falls $\widetilde{H}_i(X) \in \SC$ für alle $i \geq 0$.
\end{defn}

\begin{lem}\label{two-of-three}
  Es sei $F \to X \to B$ eine Faserung und die Räume $F$, $X$ und $B$ wegzusammenhängend.
  % TODO: Wirkung beschreiben (bisher nur für Kohomologie getan)
  Wirke $\pi_1(B)$ trivial auf $H_*(F)$.
  Dann gilt folgende $2$-aus-$3$-Eigenschaft: Falls zwei der Räume $F$, $X$ und $B$ $\SC$-azyklisch sind, so auch der dritte.
\end{lem}

\begin{proof}
  Wir betrachten die Serre-Spektralsequenz zu der Faserung mit Koeffizienten in $\Z$.
  Die Aussage, dass die Homologiegruppe $H_n(X)$ in $\SC$ liegt, ist äquivalent dazu, dass die Gruppen $E^\infty_{i,n-i}$ für $i = 0, \ldots, n$ in $\SC$ liegen, denn diese Gruppen sind die Quotienten einer endlichen Filtrierung von $H_n(X)$.

  \fall{Fall 1: $F$ und $B$ sind $\SC$-azyklisch}
  Die universelle Koeffizientenformel liefert
  \[ E^2_{pq} \cong H_p(B; H_q(F; \Z)) \cong \left( H_p(B) \otimes H_q(F) \right) \oplus \Tor(H_{p-1}(B), H_q). \]
  Man sieht durch Unterscheidung der Fälle $p=0$, $p=1$ und $p>1$ sowie $q=0$ und $q > 0$, dass $E^2_{pq} \in \SC$ für $(p, q) \neq (0, 0)$.
  Als Subquotient von $E^2_{pq}$, einer Gruppe aus $\SC$, ist dann auch $E^\infty_{pq} \in \SC$ für $(p, q) \neq (0, 0)$.
  Dies zeigt die Behauptung nach der Bemerkung am Anfang des Beweises.

  \fall{Fall 2: $F$ und $X$ sind $\SC$-azyklisch}
  Wir zeigen nun durch Induktion über $k$, dass $H_p(B) \in \SC$ für $0 < p < k$.
  %Dies gilt trivialerweise für $k = 1$.
  Gelte dies für $k \geq 1$.
  Wir wollen zeigen, dass dann auch $H_k(B)$ in $\SC$ liegt.
  %Zunächst sieht man wie im Fall 1, dass $E^2_{pq}$ und somit auch $E^r_{pq}$ in $\SC$ für $(p, q) \neq (0, 0)$ und $p < k$.
  Für alle $r \geq 2$ gibt es eine kurze exakte Sequenz
  \begin{centertikzcd}[column sep=0.4cm, row sep=0.35cm]
    0 \arrow[r] &
    \ker(d^r_{k,0}) \arrow[d, equal] \arrow[r] &
    E^r_{k,0} \arrow[r, "d^r_{k,0}"] &
    \im(d^r_{k,0}) \arrow[d, "\subseteq"] \arrow[r] &
    0 \\
    & E^{r+1}_{k,0} &&
    E^r_{k-r,r-1}
  \end{centertikzcd}
  Man sieht unter Verwendung der Induktionsannahme und der universellen Koeffizientenformel, dass $E^2_{k-r,r-1}$ und somit auch $E^r_{k-r,r-1}$ in $\SC$ liegen.
  Folglich gilt auch $\im(d^r_{k,0}) \in \SC$, also $E^r_{k,0} \in \SC \iff E^{r+1}_{k,0} \in \SC$.
  Da aber $E^R_{k,0} \cong E^\infty_{k,0} \in \SC$ für $R$ groß genug, gilt $E^r_{k,0} \in \SC$ für alle $r \geq 2$.
  Insbesondere $H_k(B; \Z) \cong H_k(B; H_0(F; \Z)) \cong E^2_{k,0} \in \SC$.

  \fall{Fall 3: $B$ und $X$ sind $\SC$-azyklisch}
  Analog zum vorherigen Fall zeigt man induktiv, dass $H_q(F) \in \SC$ für $0 < q < k$.
  Dazu verwendet man die kurze exakte Sequenz $0 \to \im(d^r_{r,k-r+1}) \hookrightarrow E^r_{0,k} \to E^{r+1}_{0,k}$.
\end{proof}

% 1.10 im Hatcher
\begin{lem}\label{homology-kgn-in-c}
  Es sei $G \in \SC \in \{ \FG, \F \}$.
  Dann ist $K(G, n)$ $\SC$-azyklisch für alle $n \geq 1$.
\end{lem}

\begin{proof}
  Sei zunächst $n=1$.
  \begin{itemize}
    \item Falls $G = \Z$, so stimmt die Aussage, denn der Kreis $S^1$ ist ein $K(\Z, 1)$ und $H_*(S^1; \Z) \in \FG$.
    \item Falls $G = \Z_m$, \TODO{begründen damit, dass der "`unendliche Linsenraum"' ein $K(\Z_m, 1)$ ist}
    \item Falls $G = G_1 \oplus G_2$, dann ist $K(G_1, 1) \times K(G_2, 1)$ ein $K(G, 1)$.
    Wenn die Aussage für $G_1$ und $G_2$ stimmt, so folgt aus dem letzten Lemma, angewendet auf die Produktfaserung $K(G_1, 1) \to K(G_1, 1) \times K(G_2, 1) \to K(G_2, 1)$, dass sie auch für $G$ gilt.
  \end{itemize}
  Da man jede endlich erzeugte abelsche Gruppe als direkte Summe von endlich vielen Summanden der Form $\Z$ und $\Z_m$ schreiben kann, gilt die Aussage für $n=1$.

  Induktiv zeigen wir nun, dass die Aussage für $n$ beliebig gilt.
  Dazu verwenden wir die Pfadraumfaserung $K(G, n) \to P \to K(G, n{+}1)$.
  Es gilt $H_k(P) = 0 \in \SC$ und $H_k(K(G, n)) \in \SC$ für $k \geq 1$ nach Induktionshypothese, also $H_k(K(G, n{+}1)) \in \SC$ für alle $k \geq 1$ nach dem vorherigen Lemma.
\end{proof}


\renewcommand\windowpagestuff{
  \vspace{1cm}
  \begin{centertikzcd}[ampersand replacement=\&, column sep=0.2cm, row sep=0.5cm]
    \&\&\& X \arrow[dll] \arrow[dl] \arrow[d] \\
    \dots \arrow[r] \&
    X_3 \arrow[r] \&
    X_2 \arrow[r] \&
    X_1
  \end{centertikzcd}
}
\opencutright
\begin{cutout}{0}{\dimexpr\linewidth-4.5cm\relax}{0pt}{2}
  \begin{defn}
    Ein \emph{Postnikov-Turm} eines wegzusammenhängenden Raumes $X$ ist ein kommutatives Diagramm wie rechts, für das gilt: \\[3pt]
    %\begin{itemize}
    %\item
    \quad$\bullet$\enspace $\pi_i(X \to X_n)$ ist ein Isomorphismus für $i \leq n$ und \\
    \quad$\bullet$\enspace $\pi_i(X_n) = 0$ für $i > n$.
    %\end{itemize}
  \end{defn}
\end{cutout}

\TODO{Bemerkung zur Konstruktion von Postnikov-Türmen}

\begin{bem}
  Es sei ein Postnikov-Turm $\ldots \to X_2 \to X_1$ gegeben.
  Dann kann man durch wiederholtes Anwenden der in Lemma~\ref{convert-to-fibration} beschriebenen Konstruktion einen neuen Postnikovturm $\ldots \to X_2' \to X_1'$ und Homotopieäquivalenzen $X_i \simeq X_i'$ konstruieren, sodass die Abbildungen $X_{i+1}' \to X_i'$ Hurewicz-Faserungen sind.
  Man sieht anhand der langen exakten Sequenz von Homotopiegruppen, dass die Faser von $X_{i+1}' \to X_i'$ ein $K(\pi_{n+1}(X), n{+}1)$ ist.
\end{bem}

\begin{lem}\label{homotopy-in-c-implies-homology-in-c}
  Es sei $\SC \in \{ \FG, \F \}$ und
  $X$ einfach zusammenhängend mit $\pi_i(X, x_0) \in \SC$ für alle $i \geq 0$.
  Dann ist $X$ $\SC$-azyklisch, \dh{} es gilt $H_i(X) \in \SC$ für $i \geq 1$.
\end{lem}

\begin{proof}
  Es sei $\ldots \to X_{i+1} \to X_i \to \ldots \to X_1$ ein Postnikov-Turm von $X$, dessen Abbildungen $X_{i+1} \to X_i$ Hurewicz-Faserungen sind.
  Wir zeigen, dass $H_i(X_k; \Z) \in \SC$ für alle $i, k > 0$.
  Die Aussage stimmt für $k = 1$, da alle Homotopiegruppen und somit auch Homologiegruppen von $X_1$ gleich Null sind.
  %Beachte, dass $X_2$ ein $K(\pi_2(X); 2)$, denn $\pi_1(X_2) \cong \pi_1(X) = 0$.
  %Da $X_2$ ist ein  ist, stimmt die Aussage nach Lemma \ref{homology-kgn-in-c} für $k=2$.
  Gelte die Aussage nun für ein $k \geq 1$.
  Wir verwenden die Faserung $K(\pi_{k+1}(X), k{+}1) \to X_{k+1} \to X_k$.
  Nach Lemma \ref{homology-kgn-in-c} sind die Homologiegruppen der Faser in $\SC$.
  Gleiches gilt für den Basisraum nach Induktionsvoraussetzung,
  und somit auch für $X_k$ nach Lemma \ref{two-of-three}.

  Es gilt $H_i(X; \Z) \cong H_i(X_k; \Z) \in \SC$ für $k \geq i$, da $\pi_i(X \to X_k)$ und nach einem Korollar der relativen Version des Hurewicz-Theorems damit auch $H_i(X \to X_k)$ ein Isomorphismus für $i \leq k$ ist.
  % TODO: Fußnote zu dem letzten Satz?
\end{proof}


\begin{proof}[Beweis von Satz \ref{hurewicz-mod-c}]
  Aufgrund der Natürlichkeit des Hurewicz-Homomorphismus kommutiert folgendes Diagramm:
  \begin{centertikzcd}
    \pi_n(X) \arrow[r, "\cong"] \arrow[d, "h_n"] &
    \pi_n(X_n) \arrow[d, "h_n"] \\
    H_n(X) \arrow[r, "\cong"] &
    H_n(X_n)
  \end{centertikzcd}
  Es genügt daher zu zeigen, dass $h_n : \pi_n(X_n) \to H_n(X_n; \Z)$ ein Isomorphismus modulo $\SC$ ist.
  Wir betrachten die Serre-Spektralsequenz zur Faserung $F_n \to X_n \to X_{n-1}$.
  \[ E^2_{pq} \cong H_p(X_{n-1}; \underbrace{H_q(K(\pi_n(X), n); \Z)}_{= 0 \text{ für } 0 < q < n}), \]
  % TODO: Begründung für das Verschwinden dieser Homologiegruppen
  Es verschwinden also alle Einträge zwischen der 0-ten und der $n$-ten Zeile. Für diese gilt
  Wir betrachten die exakte Sequenz
  \begin{centertikzcd}[row sep=0.1cm, column sep=0.5cm]
    H_{n+1}(X_{n-1}) \arrow[d, equal] &&
    H_n(F_n) \arrow[d, equal] &
    0 \arrow[rd] &&
    0 &&&
    H_n(X_{n-1}) \arrow[d, equal] \\
    E^2_{n+1,0} \arrow[d, equal] && E^2_{0,n} \arrow[d, equal] && E^\infty_{0,n} \arrow[rrd] \arrow[ru] &&&&
    E^2_{n,0} \arrow[d, equal] \\
    %\underbrace{E^n_{n+1,0}}_{\mathrlap{\substack{= E^2_{n+1,0} \\ = H_{n+1}(X_{n-1})}}} \arrow[r, "d^n_{n+1,0}"] &
    E^n_{n+1,0} \arrow[rr, "d^r_{n+1,0}"] &&
    E^n_{0,n} \arrow[rrrr] \arrow[rru] &&&&
    H_n(X_n) \arrow[rr] &&
    E^\infty_{n,0} \arrow[r] &
    0
  \end{centertikzcd}
  welche sich aus zwei kurzen Sequenzen zusammensetzt:
  \begin{itemize}
    \item Die linke Sequenz ist exakt, da $E^\infty_{0,n} \cong E^{n+1}_{0,n} \cong \ker(d^r_{0,n}) / \im(d^r_{n+1,0}) = E^n_{0,n} / \im(d^r_{n+1,0})$.
    \item Die rechte Sequenz ist exakt, da $E^\infty_{0,n} = F^0 H_n(X_n)$ und $E^\infty_{n,0} = F^n H_n(X_n) / F^{n-1} H_n(X_n)$ für eine Filtrierung $0 = F^{-1} H_n(X_n) \subseteq F^0 H_n(X_n) \subseteq \ldots \subseteq F^n H_n(X_n) = H_n(X_n)$.
    Da $F^p H_n(X_n) / F^{p-1} H_n(X_n) = E^\infty_{p,n-p} = 0$ für $p = 1, \ldots, n-1$, gilt $F^0 H_n(X_n) = \ldots = F^{n-1} H_n(X_n)$.
  \end{itemize}
  Aus Lemma \label{homotopy-in-c-implies-homology-in-c} folgt, dass $E^n_{n+1,0}, E^\infty_{n,0} \in \SC$.
  Somit ist der mittlere Morphismus $H_n(F_n) = E^n_{0,n} \to H_n(X_n)$ ein Isomorphismus modulo $\SC$.
  \TODO{Begründen, warum dieser Morphismus genau der von der Inklusion $F_n \hookrightarrow X_n$ induzierte Morphismus ist.}
  Betrachte nun das kommutative Diagramm
  \begin{centertikzcd}
    \pi_n(F_n) \arrow{d}{h_n}[swap]{\cong} \arrow[r, "\cong"] &
    \pi_n(X_n) \arrow[d, "h_n"] \\
    H_n(F_n) \arrow[r] &
    H_n(X_n)
  \end{centertikzcd}
  Dabei sind die vertikalen Morphismen die Hurewicz-Homomorphismen und die horizontalen Morphismen werden durch die Inklusion induziert.
  Der linke Hurewicz-Homomorphismus ist nach dem Hurewicz-Theorem ein Isomorphismus, da $F_n$ $(n{-}1)$-zusammenhängend ist.
  Der obere Morphismus ist ein Isomorphismus, wie man anhand der langen exakten Sequenz der Faserung $F_n \to X_n \to X_{n-1}$ sehen kann.
  Der untere Morphismus ist ein Isomorphismus modulo $\SC$, wie wir gerade eben gesehen haben.
  Somit ist auch der rechte Morphismus im Diagramm ein Isomorphismus modulo $\SC$.
\end{proof}

\begin{bem}
  Man kann in diesem Kapitel die Voraussetzung, dass $X$ einfach zusammenhängend ist, ersetzen durch die Forderung, dass $X$ wegzusammenhängend und abelsch ist, \dh{} die Wirkung der Fundamentalgruppe $\pi_1(X)$ auf den höheren Homotopiegruppen $\pi_n(X)$ trivial ist.
\end{bem}

\section{Rationale Kohomologie von Räumen vom Typ $K(\Z, n)$}

% TODO: Etwas schreiben über rationale Homotopiegruppen

% Hatcher, 1.20
\begin{satz}\label{rational-homology-kzn}
  Für $n \geq 1$ gilt
  \[
    H^*(K(\Z, n); \Q) \cong \begin{cases}
      \Q[x], & \text{falls $n$ gerade}, \\
      \Lambda_\Q[x], & \text{falls $n$ ungerade}
    \end{cases}
  \]
  als graduierte Ringe mit Erzeuger $x \in H^n(K(\Z, n); \Q)$.
  Dabei bezeichnet $\Lambda_\Q[x]$ die äußere Algebra mit Erzeuger $x$.
\end{satz}

\begin{proof}
  Durch Induktion über $n$.
  Der Satz gilt für $n = 1$, denn der Kreis $S^1$ ist ein $K(\Z, 1)$ und es gilt bekanntermaßen $H^*(S^1; R) \cong \Lambda_R[x]$ für $R = \Z$ und somit auch für $R = \Q$. \\
  Im Induktionsschritt nutzen wir die Pfadraumfaserung $F \coloneqq K(\Z, n{-}1) \to P \to B \coloneqq K(\Z, n)$. Da $K(\Z, n)$ für $n \geq 2$ einfach zusammenhängend ist, gilt für deren zugehörige Serre-Spektralsequenz
  $E_2^{p,q} \cong H^p(B; H^q(F))$.

  \vspace{0.5cm}

  \renewcommand\windowpagestuff{
    \begin{center}\begin{tikzpicture}[x=16pt,y=16pt]\begin{scope}[shift={(0,0)}]
      \foreach \x in {1,2,3,5,6,7}{
        \foreach \y in {0,1,...,3}{
          \zeroDot{\x}{\y}
        }
      }
      \foreach \x in {0,...,8}{
        \zeroDot{\x}{4}
      }
      \foreach \y in {0,...,4}{
        \node at (9.5,\y) {$\cdots$};
      }
      \node at (0,0) {$\Q$\small $1$};
      \zeroDot{0}{1}
      \zeroDot{0}{2}
      \node at (0,3) {$\Q$\small $a$};
      \node at (4,0) {$\Q$\small $x$};
      \zeroDot{4}{1}
      \zeroDot{4}{2}
      \node at (4,3) {$\Q$\small $ax$};
      \draw[->,gray] (0.6,2.6) -- (3.4,0.4);
      \node at (8,0) {$\Q$\small $x^2$};
      \zeroDot{8}{1}
      \zeroDot{8}{2}
      \node at (8,3) {$\Q$\small $ax^2$};
      \draw[->,gray] (4.6,2.6) -- (7.4,0.4);
      \draw[->] (-0.7,-0.6) -- (10.3,-0.6); %node[below] {p};
      \draw[->] (-0.7,-0.6) -- (-0.7,4.5); %node[left] {q};
      \node[left] at (-0.7,0) {$0$};
      \node[left] at (-0.9,1.7) {\vdots};
      \node[left] at (-0.7,3) {$n{-}1$};
      \node[below] at (0,-0.6) {$0$};
      \node[below] at (2,-0.7) {$\cdots$};
      \node[below] at (4,-0.7) {$n$};
      \node[below] at (6,-0.7) {$\cdots$};
      \node[below] at (8,-0.6) {$2n$};
      \node[below] at (9.5,-0.7) {$\cdots$};
    \end{scope}\end{tikzpicture}\end{center}
  }
  \opencutright
  \begin{cutout}{0}{\dimexpr\linewidth-8cm\relax}{0pt}{8}
    \emph{Falls $n$ gerade:} \enspace
    Dann sieht die Spektralsequenz auf der Seite $E_r$, $r \leq n$ aus wie rechts skizziert (dabei stehen die Gitterpunkte für die Nullgruppe).
    Das sieht man folgendermaßen: Zunächst ist $E_2^{pq} = 0$ und somit $E_r^{pq} = 0$ außer für $q \in \{ 0, n{-}1 \}$, denn nach Induktionsvoraussetzung gilt $H^*(F; \Q) \cong \Lambda_\Q[x]$.
    Es folgt, dass nur auf der $n$-ten Seite $E_n$ nicht verschwindende Differentiale existieren können und $E_2 \cong E_n$ und $E_{n+1} \cong E_\infty$ gilt.
    Außerdem ist $E_n^{0,0} \cong E_2^{0,0} \cong \Q$ und $E_n^{0,n-1} \cong E_2^{0,n-1} \cong \Q$, da $B$ zusammenhängend ist.
    Die Spektralsequenz konvergiert gegen $H^*(P; \Q)$.
    Da $P$ zusammenziehbar ist, gilt $H^0(P; \Q) = \Q$ und $H^n(P; \Q) = 0$ für $n > 0$.
    Folglich ist $E_{n+1}^{pq} \cong E_\infty^{pq} = 0$ außer für $p = q = 0$.
    Insbesondere gilt $E_{n+1}^{0,n-1} \cong E_\infty^{0,n-1} = 0$.
    %Es gilt $E_n^{0,n-1} \cong E_2^{0,n-1} \cong \Q$, da die von $E_r^{0,n-1}$ ausgehenden Differentiale $d_r^{0,n-1}$ für $r < n$ Null und die dort eintreffenden Differentiale für alle $r$ Null sind.
    Das eingezeichnete Differential $d_n^{0,n-1} : E_n^{0,n-1} \to E_n^{n,0}$ ist nun injektiv, denn $\ker(d_n^{0,n-1}) \cong E_{n+1}^{0,n-1} = 0$.
    Dieses Differential ist auch surjektiv, denn $\coker(d_n^{0,n-1}) \cong E_{n+1}^{n,0} \cong E_\infty^{n,0} = 0$, also ein Isomorphismus.
    Somit $H^n(B; \Q) \cong E_2^{n,0} \cong E_n^{n,0} \cong E_n^{0,n-1} \cong \Q$.
    Der zweite Isomorphismus kommt daher, da für $r \leq n-1$ alle Differentiale von und nach $E_r^{n,0}$ Null sind. % XXX: besser formulieren?
    Damit ist $E_2^{n,n-1} \cong H^{n}(B; H^{n-1}(F; \Q)) \cong H^{n}(B; \Q) \cong \Q$.
    Induktiv sieht man nun, dass die Abbildungen $d_n^{kn,n-1}$ Isomorphismen sind und dass $H^{kn}(B; \Q) \cong E_2^{kn,0} \cong  E_2^{kn,n-1} \cong \Q$ für alle $k \geq 0$.
    Damit haben wir gezeigt, dass die graduierte, additive Struktur von $H^*(B; \Q)$ wie behauptet ist.
  \end{cutout}
  Es sei nun $a \in E_n^{0,n-1} \cong H^0(B; H^{n-1}(F; \Q))$ ungleich Null und $x \coloneqq d_n^{0,n-1}(a) \in E_n^{n,0} \cong H^n(B; H^0(F; \Q)) \cong H^n(B; \Q)$.
  Dann gilt auch $x \neq 0$ und $ax \coloneqq m_r(a, x) \neq 0$, da wegen (ii) und (iii) das Produkt $m_r$ gerade dem kanonischen Produkt $H^n(B; H^0(F; \Q)) \times H^0(B; H^{n-1}(F; \Q)) \to H^n(B; H^{n-1}(F; \Q))$ entspricht. % XXX: kommutatives Diagramm dazu?
  Es gilt
  $0 \neq d_n^{n,n-1}(ax) = d_n^{0,n-1}(a)x - a d_n^{n,0}(x) = xx$.
  Da das Produkt $xx \in E_n^{2n,0}$ gerade dem Cup-Produkt $x \cup x \in H^{2n}(B; \Q)$ entspricht, ist $x \cup x \neq 0$, also ein Erzeuger von $H^{2n}(B; \Q)$.
  Induktiv ist nun $0 \neq d_n^{kn,n-1}(a x^k) = x^{k+1} \in E_n^{kn,0}$ da ja $0 \neq a x^k$.
  Somit ist für alle $k$ das $k$-fache Cup-Produkt $x^k \in H^{kn}(B; \Q)$ ein Erzeuger.

  \vspace{1.5cm}

  \renewcommand\windowpagestuff{
    \begin{center}\begin{tikzpicture}[x=16pt,y=16pt]\begin{scope}[shift={(0,0)}]
      \foreach \x in {1,2,3}{
        \foreach \y in {0,1,...,6}{
          \zeroDot{\x}{\y}
        }
      }
      \foreach \y in {0,...,6}{
        \zeroDot{5}{\y}
      }
      \foreach \x in {0,...,5}{
        \node at (\x,7.3) {\vdots};
      }
      \node at (0,0) {$\Q$\small $1$};
      \zeroDot{0}{1}
      \zeroDot{0}{2}
      \node at (0,3) {$\Q$\small $a$};
      \zeroDot{0}{4}
      \zeroDot{0}{5}
      \node at (0,6) {$\Q$\small $a^2$};
      \node at (4,0) {$\Q$\small $x$};
      \zeroDot{4}{1}
      \zeroDot{4}{2}
      \node at (4,3) {$\Q$\small $ax$};
      \zeroDot{4}{4}
      \zeroDot{4}{5}
      \node at (4,6) {$\Q$\small $a^2 x$};
      \draw[->,gray] (0.6,2.6) -- (3.4,0.4);
      \draw[->,gray] (0.6,5.6) -- (3.4,3.4);
      \draw[->] (-0.7,-0.6) -- (5.6,-0.6); %node[below] {p};
      \draw[->] (-0.7,-0.6) -- (-0.7,7.8); %node[left] {q};
      \node[left] at (-0.7,0) {$0$};
      \node[left] at (-0.9,1.7) {\vdots};
      \node[left] at (-0.7,3) {$n{-}1$};
      \node[left] at (-0.9,4.7) {\vdots};
      \node[left] at (-0.7,6) {$2n{-}2$};
      \node[below] at (0,-0.6) {$0$};
      \node[below] at (2,-0.7) {$\cdots$};
      \node[below] at (4,-0.7) {$n$};
    \end{scope}\end{tikzpicture}\end{center}
  }
  \opencutleft
  \begin{cutout}{0}{0pt}{\dimexpr\linewidth-5.5cm\relax}{11}
    \emph{Falls $n$ ungerade:} \enspace
    Dann ist $E_r^{p,q} = 0$ für alle $q$, die kein Vielfaches von $n-1$ sind.
    Somit verschwinden alle Differentiale auf $E_r$ für $r < n$ und $E_2 \cong E_n$.
    Für $0 < m < n$ verschwinden alle Differentiale von und nach $E_r^{m,0}$ und daher ist $H^m(B; \Q) \cong E_2^{m,0} \cong E_\infty^{m,0} = 0$ und  folglich $E_2^{m,k} = 0$ für alle $k \geq 0$.
    Selbiges gilt folglich auch für $n < m < 2n$ und allgemeiner für solche $m$, die kein Vielfaches von $n$ sind.
    Analog wie im vorherigen Fall sieht man, dass das eingezeichnete Differential $d_n^{0,n-1}$ ein Isomorphismus ist.
    Somit $H^n(B; \Q) \cong H^n(B; H^0(F; \Q)) \cong \Q$ und $E_2^{n,k(n-1)} \cong \Q$ für alle $k \geq 0$.
    Sei $a \in H^0(B; H^{n-1}(F; \Q))$ ungleich Null und $x \coloneqq d_n^{0,n-1}$.
    Dann ist auch $a^2 \neq 0 \in E_n^{0,2n-2}$ und $d_n^{0,2n-2}(a^2) = d_n^{0,n-1}(a)a + d_n^{0,n-1}(a)a = xa + ax = (-1)^{0 \cdot n + (n-1) \cdot 0} ax + ax = 2ax \neq 0$.
    Also ist $d_n^{0,2n-2}$ ein Isomorphismus.
    Analog sieht man, dass $d_n^{0,k(n-1)}$ für alle $k \geq 1$ ein Isomorphismus ist.
    Es bleibt zu zeigen, dass $H^{kn}(B; \Q) = 0$ für $k > 1$.
    Das einzige potentiell nichttriviale Differential, das bei $E_r^{2n,0}$ ankommt, ist $d_n^{n,n-1}$.
    Dieses ist aber Null, da $\ker(d_n^{n,n-1}) = \im(d_n^{0,2n-2}) = E_n^{n,n-1}$.
    Also $H^{2n}(B; \Q) \cong E_2^{2n,0} \cong E_\infty^{2n,0} = 0$ und $E_2^{2n,k} = 0$ für alle $k \geq 0$.
    Für $k > 2$ sieht man durch Induktion, dass alle Differentiale von und nach $E_r^{2n,0}$ verschwinden und daher $H^{kn}(B; \Q) = 0$.
  \end{cutout}
\end{proof}

\section{Beweis des Satzes von Serre}

\begin{lem}
  Es sei $n \geq 3$ ungerade und $X$ ein $(n{-}1)$-zusammenhängender topologischer Raum.
  Angenommen, $H^k(X; \Z) = 0$ für $k > n$ und $H_n(X) \cong \pi_n(X)$ ist die direkte Summe von $\Z$ und einer endlichen Gruppe.
  Dann sind die Homotopiegruppen $\pi_k(X, x_0)$, $k > n$ endlich.
\end{lem}

\begin{proof}
  Durch Töten der höheren Homotopiegruppen bekommen wir eine Abbildung $f : X \to K(\pi_n(X), n) \approx K(\Z, n) \times K(G, n)$, wobei $G$ eine endliche Gruppe ist.
  Wir führen die in Lemma~\ref{convert-to-fibration} beschriebene Konstruktion durch und erhalten einen zu $X$ homotopieäquivalenten Raum~$X'$ und eine Hurewicz-Faserung $p : X' \to K(\pi_n(X), n)$ mit Faser $F$.
  Anhand der langen exakten Sequenz dieser Faserung sehen wir, dass
  \[
    \pi_i(F) \cong \begin{cases}
      0, & \text{für $i \leq n$}, \\
      \pi_i(X') \cong \pi_i(X), & \text{für $i > n$.}
    \end{cases}
  \]
  Wir wenden diesselbe Konstruktion auf die Inklusion $F \hookrightarrow X'$ an und bekommen einen zu $F$ homotopieäquivalenten Raum $F'$ und eine Faserung $q : F' \to X'$.
  Aus der langen exakten Homotopiesequenz ergibt sich, dass die Faser~$\tilde{F}$ dieser Faserung ein $K(\pi_n(X), n{-}1)$ ist.

  Wir wollen die Serre-Spektralsequenz der Faserung $\tilde{F} \to F' \to X'$ mit Koeffizienten in~$\Q$ verwenden.
  Dazu untersuchen wir zunächst die rationale Kohomologie von Faser und Basisraum.
  Lemma \ref{homology-kgn-in-c} impliziert, dass die Homotopiegruppen und wegen Korollar \ref{homotopy-in-c-iff-homology-in-c} auch die reduzierten Homologiegruppen von $K(G, n{-}1)$ endlich sind.
  Nach der universellen Koeffizientenformel verschwinden somit alle reduzierten Kohomologiegruppen von $K(G, n{-}1)$ mit Koeffizienten in~$\Q$.
  Wir sehen nun an der Serre-Spektralsequenz der Produktfaserung $K(G, n{-}1) \to K(\pi_n(X), n{-}1) \to K(\Z, n{-}1)$ und Lemma \ref{rational-homology-kzn}, dass $H^*(K(\pi_n(X), n{-}1); \Q) \cong H^*(K(\Z, n{-}1); \Q) \cong \Q[x]$ mit Erzeuger $x \in H^{n-1}(K(\pi_n(X), n{-}1); \Q)$.
  Für den Basisraum folgt aus der universellen Koeffizientenformel, dass die rationale Kohomologie von~$X'$ gleich~$\Q$ ist in Grad~$0$ und in Grad~$n$, und Null sonst.

  Wir wissen nun, dass $E_2^{pq} = 0$ außer falls $p \in \{ 0, n \}$ und $n{-}1 \divides q$ gilt.
  Die Spektralsequenz besitzt also auf der $E_2$-Seite und damit auch auf der $E_n$-Seite die gleichen Einträge wie die Spektralsequenz aus dem zweiten Fall ($n$ ungerade) des vorhergehenden Lemmas.
  Genau wie dort schließen wir, dass $d_n^{0,n-1}$ ein Isomorphismus ist (denn $H^n(X'; \Q) = 0$) und dass aufgrund der multiplikativen Struktur der Spektralsequenz auch die Differentiale $d_n^{0,k(n-1)}$ für $k \geq 1$ Isomorphismen sind.
  Somit ist $E_{n+1}^{pq} = 0$ außer für $p = q = 0$.
  Es folgt, dass $H^n(F'; \Q) = 0$ für $n > 0$.
  Somit gilt auch $H_n(F'; \Z) \times \Q \cong H_n(F'; \Q) \cong H^n(F'; \Q) = 0$, es ist also $H_n(F'; \Z)$ eine Torsionsgruppe.
  Da $X' \simeq X$ und $K(\pi_n(X), n)$ $\FG$-azyklisch sind, ist auch $F' \simeq F$ $\FG$-azyklisch, \dh{} die Homologiegruppen $H_n(F'; \Z)$,~$n > 0$ sind endlich erzeugt.
  Somit sind diese Homologiegruppen schon endlich.
  Nach Korollar~\ref{homotopy-in-c-iff-homology-in-c} sind damit alle Homotopiegruppen von $F'$ endlich.
  Diese entsprechen im Grad $k > n$ aber gerade den Homotopiegruppen von $X$.
\end{proof}

\begin{proof}[Beweis des Satzes von Serre (Satz \ref{})]
  Im Fall $n=1$ stimmt die Aussage, denn die universelle Überlagerung von $S^1$ ist $\R$.
  Wir können daher $n \geq 2$ annehmen.
  Im Fall, dass $n \geq 3$ ungerade ist, folgt die Aussage aus dem vorhergehenden Lemma.
  Es bleibt der Fall, dass $n \geq 2$ gerade ist.

  \vspace{2em}

  \renewcommand\windowpagestuff{
    \begin{center}\begin{tikzpicture}[x=16pt,y=16pt]\begin{scope}[shift={(0,0)}]
      \foreach \x in {1,2,3,5}{
        \foreach \y in {0,1,...,3}{
          \zeroDot{\x}{\y}
        }
      }
      \foreach \x in {0,...,5}{
        \zeroDot{\x}{4}
      }
      \node at (0,0) {$\Q$\small $1$};
      \zeroDot{0}{1}
      \zeroDot{0}{2}
      \node at (0,3) {$\Q$\small $a$};
      \node at (4,0) {$\Q$\small $x$};
      \zeroDot{4}{1}
      \zeroDot{4}{2}
      \node at (4,3) {$\Q$\small $ax$};
      \draw[->,gray] (0.6,2.6) -- (3.4,0.4);
      \draw[->] (-0.7,-0.6) -- (6,-0.6); %node[below] {p};
      \draw[->] (-0.7,-0.6) -- (-0.7,4.5); %node[left] {q};
      \node[left] at (-0.7,0) {$0$};
      \node[left] at (-0.9,1.7) {\vdots};
      \node[left] at (-0.7,3) {$n{-}1$};
      \node[below] at (0,-0.6) {$0$};
      \node[below] at (2,-0.7) {$\cdots$};
      \node[below] at (4,-0.7) {$n$};
    \end{scope}\end{tikzpicture}\end{center}
  }
  \opencutright
  \begin{cutout}{0}{\dimexpr\linewidth-5cm\relax}{0pt}{8}
    Wir wenden diesselbe Konstruktion wie im Beweis des vorhergehenden Lemmas für den Raum $X = S^n$ an.
    Da $n$ gerade ist, haben wir aber $H^*(K(\pi_n(X), n{-}1)) \cong H^*(K(\pi_n(X), n{-}1)) \cong \Lambda[x]$.
    %Also $E_2^{pq} = 0$ außer für $(p, q) \in \{ 0, n \} \times \{ 0, n-1 \}$.
    Die $E_n$-Seite der Serre-Spektralsequenz der Faserung $K(\pi_n(X), n{-}1) \to F' \to X'$ ist daher wie rechts abgebildet.
    Der Eintrag $\Q$ auf Position $E_2^{n,n-1}$ überlebt die Spektralsequenz.
    Somit ist $H^*(F'; \Q) \cong \Lambda[b]$ mit $b \in H^{2n-1}(F'; \Q)$.
    Es folgt $H_i(F'; \Z) \otimes \Q \cong H_i(F'; \Q) = 0$ für $i < 2n-1$.
    Da die Homologiegruppen von $F'$ endlich erzeugt sind, ist $H_i(F'; \Z) \in \F$ für $i < 2n-1$ und $H_{2n-1}(F'; \Z) \cong \Z \oplus G$, wobei $G$ eine endliche Gruppe ist.
    Das Hurewicz-modulo-$\SC$-Theorem impliziert, dass damit auch die Homotopiegruppen $\pi_i(F')$ endlich sind für $i < 2n-1$ und dass $\pi_{2n-1}(F) \cong \Z \oplus G'$ für eine endliche Gruppe $G'$.
    Indem wir höhere Homotopiegruppen killen, erhalten wir eine Abbildung $F' \to Y$, wobei $\pi_i(Y) = 0$ für $i \geq 2n-1$.
    Wir konvertieren diese Abbildung zu einer Faserung $F'' \to Y$ mit Faser $Z$.
    Anhand der langen exakten Homotopiesequenz dieser Faserung sehen wir, dass
    \[
      \pi_m(Z) = \begin{cases}
        0, & \text{für $m < 2n-1$,} \\
        \pi_m(F'') \cong \pi_m(F) \cong \pi_m(S^n), & \text{für $m \geq 2n-1$.}
      \end{cases}
    \]
    Wir bemerken, dass alle Homotopiegruppen und somit auch Homologiegruppen von~$Y$ endlich sind.
    Daher gilt $\widetilde{H}^*(Y; \Q) = 0$.
    Da die Serre-Spektralsequenz zur Faserung $Z \to F'' \to Y$ mit Koeffizienten in~$\Q$ auf der $E_2$-Seite also nur Einträge in der Spalte $p=0$ besitzt, gilt $H^*(Z; \Q) \cong H^*(F''; \Q) \cong \Lambda[b]$.
    Die Aussage folgt nun aus dem vorhergehenden Lemma mit $X \coloneqq Z$.
  \end{cutout}
\end{proof}

\nocite{*}
\addcontentsline{toc}{section}{Literatur}
\printbibliography

\newpage
\addcontentsline{toc}{section}{Anhang: Tabelle der Homotopiegruppen der Sphären}
\begin{landscape}
  \begin{center}
    \mbox{}\vspace{8em}
    \begin{table}
      \footnotesize \setlength\extrarowheight{5pt}
      \begin{tabular}{r | c c c c c c c c c c c | l}
  $\pi_{n+k}(S^n)$ & $n{=}1$ & $n{=}2$ & $n{=}3$ & $n{=}4$ & $n{=}5$ & $n{=}6$ & $n{=}7$ & $n{=}8$ & $n{=}9$ & $n{=}10$ & $n{=}11$ & $\pi^s_k$ \\[0.2em] \hline
  %$k < 0$ & $0$ & $\cdots$ \\[2pt] \hhline{~-}
  $k = 0$ & \multicolumn{1}{c|}{\cellcolor{gray!15} $\Z$} & \cellcolor{gray!15} $\Z$ & \cellcolor{gray!15} $\cdots$ & \cellcolor{gray!15} $\cdots$ & \cellcolor{gray!15} $\cdots$ & \cellcolor{gray!15} $\cdots$ & \cellcolor{gray!15} $\cdots$ & \cellcolor{gray!15} $\cdots$ & \cellcolor{gray!15} $\cdots$ & \cellcolor{gray!15} $\cdots$ & \cellcolor{gray!15} $\cdots$ & \cellcolor{gray!15} $\Z$ \\[2pt] \hhline{~~-}
  $k = 1$ & $0$ & \multicolumn{1}{c|}{\cellcolor{gray!15} $\Z$} & $\ZZ{2}{}$ & $\cdots$ & $\cdots$ & $\cdots$ & $\cdots$ & $\cdots$ & $\cdots$ & $\cdots$ & $\cdots$ & $\ZZ{2}{}$ \\[2pt] \hhline{~~~-}
  $k = 2$ & $0$ & $\ZZ{2}{}$ & \multicolumn{1}{c|}{$\ZZ{2}{}$} & $\ZZ{2}{}$ & $\cdots$ & $\cdots$ & $\cdots$ & $\cdots$ & $\cdots$ & $\cdots$ & $\cdots$ & $\ZZ{2}{}$ \\[2pt] \hhline{~~~~-}
  $k = 3$ & $0$ & $\ZZ{2}{}$ & $\ZZ{4}{} {\times} \ZZ{3}{}$ & \multicolumn{1}{c|}{\cellcolor{gray!15} $\Z {\times} \ZZ{4}{} {\times} \ZZ{3}{}$} & $\ZZ{8}{} {\times} \ZZ{3}{}$ & $\cdots$ & $\cdots$ & $\cdots$ & $\cdots$ & $\cdots$ & $\cdots$ & $\ZZ{8}{} {\times} \ZZ{3}{}$ \\[2pt] \hhline{~~~~~-}
  $k = 4$ & $0$ & $\ZZ{4}{} {\times} \ZZ{3}{}$ & $\ZZ{2}{}$ & $\ZZ{2}{2}$ & \multicolumn{1}{c|}{$\ZZ{2}{}$} & $0$ & $\cdots$ & $\cdots$ & $\cdots$ & $\cdots$ & $\cdots$ & $0$ \\[2pt] \hhline{~~~~~~-}
  $k = 5$ & $0$ & $\ZZ{2}{}$ & $\ZZ{2}{}$ & $\ZZ{2}{2}$ & $\ZZ{2}{}$ & \multicolumn{1}{c|}{\cellcolor{gray!15} $\Z$} & $0$ & $\cdots$ & $\cdots$ & $\cdots$ & $\cdots$ & $0$ \\[2pt] \hhline{~~~~~~~-}
  $k = 6$ & $0$ & $\ZZ{2}{}$ & $\ZZ{3}{}$ & $\ZZ{8}{} {\times} \ZZ{3}{2}$ & $\ZZ{2}{}$ & $\ZZ{2}{}$ & \multicolumn{1}{c|}{$\ZZ{2}{}$} & $\ZZ{2}{}$ & $\cdots$ & $\cdots$ & $\cdots$ & $\ZZ{2}{}$ \\[2pt] \hhline{~~~~~~~~-}
  $k = 7$ & $0$ & $\ZZ{3}{}$ & $\ZZ{3}{} {\times} \ZZ{5}{}$ & $\ZZ{3}{} {\times} \ZZ{5}{}$ & $\ZZ{2}{} {\times} \ZZ{3}{} {\times} \ZZ{5}{}$ & $\ZZ{4}{} {\times} \ZZ{3}{} {\times} \ZZ{5}{}$ & $\ZZ{8}{} {\times} \ZZ{3}{} {\times} \ZZ{5}{}$ & \multicolumn{1}{c|}{\cellcolor{gray!15} $\Z {\times} \ZZ{8}{} {\times} \ZZ{3}{} {\times} \ZZ{5}{}$} & $\ZZ{16}{} {\times} \ZZ{3}{} {\times} \ZZ{5}{}$ & $\cdots$ & $\cdots$ & $\ZZ{16}{} {\times} \ZZ{3}{} {\times} \ZZ{5}{}$ \\[2pt] \hhline{~~~~~~~~~-}
  $k = 8$ & $0$ & $\ZZ{3}{} {\times} \ZZ{5}{}$ & $\ZZ{2}{}$ & $\ZZ{2}{}$ & $\ZZ{2}{}$ & $\ZZ{8}{} {\times} \ZZ{2}{} {\times} \ZZ{3}{}$ & $\ZZ{2}{3}$ & $\ZZ{2}{4}$ & \multicolumn{1}{c|}{$\ZZ{2}{3}$} & $\ZZ{2}{2}$ & $\cdots$ & $\ZZ{2}{2}$ \\[2pt] \hhline{~~~~~~~~~~-}
  $k = 9$ & $0$ & $\ZZ{2}{}$ & $\ZZ{2}{2}$ & $\ZZ{2}{3}$ & $\ZZ{2}{3}$ & $\ZZ{2}{3}$ & $\ZZ{2}{4}$ & $\ZZ{2}{5}$ & $\ZZ{2}{4}$ & \multicolumn{1}{c|}{\cellcolor{gray!15} $\Z {\times} \ZZ{2}{3}$} & $\ZZ{2}{3}$ & $\ZZ{2}{3}$ \\[2pt]
  %$k = 10$ & $0$ & $\ZZ{2}{2}$ & $\ZZ{4}{} {\times} \ZZ{2}{} {\times} \ZZ{3}{}$ & $\ZZ{8}{} {\times} \ZZ{4}{} {\times} \ZZ{2}{} {\times} \ZZ{3}{2} {\times} \ZZ{5}{}$ & $\ZZ{8}{} {\times} \ZZ{2}{} {\times} \ZZ{9}{}$ & $\ZZ{8}{} {\times} \ZZ{2}{} {\times} \ZZ{9}{}$ & $\ZZ{8}{} {\times} \ZZ{3}{} {\times} \ZZ{2}{}$ & $\ZZ{8}{2} {\times} \ZZ{2}{} {\times} \ZZ{3}{2}$ & $\ZZ{8}{} {\times} \ZZ{3}{} {\times} \ZZ{2}{}$ & $\ZZ{4}{} {\times} \ZZ{3}{} {\times} \ZZ{2}{}$ & $\ZZ{2}{2} {\times} \ZZ{3}{}$ & $\ZZ{2}{} {\oplus} \ZZ{3}{}$
\end{tabular}

      \caption{Die ersten Homotopiegruppen der Sphären. Der stabile Bereich besteht aus allen Gruppen oberhalb der Treppenlinie. Unendliche Gruppen sind grau hinterlegt. Quelle: \cite{toda:composition-methods}}
    \end{table}
  \end{center}
\end{landscape}

% TODO für 31.7.2015:
% * Erklären, woher man pi_i(S^n) für i <= n kennt (Einhängungssatz, simpliziale Approximation)
% * Technik des Killens von Homotopiegruppen und Konstruktion von Eilenberg-MacLane-Räumen erklären
% * TODOs wegmachen

% TODO:
% * Beweis der Existenz der Serre-Spektralsequenz
% * Eilenberg-MacLane-Räume
% simpliziale Approximation
% pi_q(S^n) wobei q < n
% pi_n(S^n) = Z (mit Freudenthal-Suspensionssatz?)
% pi_n(S^1)
% * Satz von Serre (1.21)
% * Was hat die Hopf-Invariante mit dem ganzen Zeugs zu tun?
% * Existenz der Homologie-Spektralsequenz
% * Mehr Vorsicht bei Voraussetzungen an Existenz der Serre-Spektralsequenz: Z.B. ist Raum einfach zusammenhängend?
% * Methode des Killens von unteren Homotopiegruppen separat beschreiben (und damit Beweise vereinfachen?)
% Beweise umformulieren: "man" vermeiden

% Theoreme, die benutzt werden, die man beweisen müsste:
% * Zelluläre Approximation

\end{document}
